%============================================================
\part{理論}\label{part:theory}
%============================================================

\chapter{序論：現金の時間構造を変数に選ぶ}

\section{問題設定}

企業を記述する標準的な変数は、売上高、利益率、市場占有率である。
本稿はこれらを用いず、\textbf{現金の時間構造}からビジネスモデルを記述する。

この選択には理由がある。企業の存続条件は損益ではなく現金残高に対する制約として与えられ、
しかもその制約は期待値ではなく経路ごとに課される（第\ref{sec:objective}節）。
存続と成長を扱う枠組みは、現金の時間構造を第一級の対象としなければならない。

標準的な変数を捨てるわけではない。
\textbf{損益計算書、貸借対照表、キャッシュフロー計算書は、
本稿が構成する対象を添字方向に潰した射影として現れる}（表\ref{tab:statements-projection}）。
本稿は会計と別の量を置くのではなく、会計が潰している添字を残す。

\section{本稿の構成と読み方}

本稿は六部からなる。

\begin{center}
\begin{tabularx}{\textwidth}{@{}llX@{}}
\toprule
部 & 表題 & 内容 \\
\midrule
\ref{part:theory} & 理論 & 定義、命題、証明。以降のすべての基礎 \\
\ref{part:catalog} & 写像の空間 & $\Phi$ の全空間（28類型）、制約下の部分空間、その極限、情報構造を差し替えた場合 \\
\ref{part:method} & 方法論 & 観測に接続する五段階、選択バイアスの扱い、検証障害の分類 \\
\ref{part:results} & 実証 & 理論量の測定対応と四領域の検証 \\
\ref{part:conclusion} & 総括 & 到達点と未解決問題 \\
\ref{app:proofs} & 付録 & 方法論の証明、記号一覧、既存理論との対応、データ源、事前登録、改訂履歴 \\
\bottomrule
\end{tabularx}
\captionof{table}{本稿の部構成。}
\label{tab:part-structure}
\end{center}

第\ref{part:theory}部と第\ref{part:catalog}部は演繹的であり、
第\ref{part:method}部は理論と実証を接続する手続きを与える。

第\ref{part:catalog}部は三つの段階からなる。
まず $\Phi$ の自由度から生成される全空間を列挙し（第\ref{ch:catalog}章）、
次に資本と人手の制約を課したときの部分空間を求め（第\ref{ch:solo}章）、
最後に運営者を除去する極限を扱う（第\ref{part:unmanned}章）。
いずれも演繹であり、実証ではない。

独立した「応用」部を設けなかったのは意図的である。
応用が正当化されるのは理論が実証を経た後であり、本稿の実証は完結していない。
制約下の演繹は応用ではなく、\textbf{写像の空間を制約で切り出す操作}として
第\ref{part:catalog}部に置いた。

第\ref{part:results}部の冒頭に第\ref{ch:instantiation}章を置き、
理論で導入した記号のそれぞれについて、
実測できた値と測定できなかった事実を一覧した。
以降の各章にも冒頭に「本章で扱う理論量」の表を付し、
どの記号がどこで測られたかを追えるようにしてある。

読者の関心に応じて、次の経路を推奨する。

\begin{itemize}
\item \textbf{理論のみに関心がある場合}：第\ref{part:theory}部と第\ref{part:catalog}部。
命題の証明は本文に置いてある。方法論上の主張の証明のみ付録\ref{app:proofs}にまとめた
\item \textbf{実務的な設計に関心がある場合}：第\ref{sec:growth}章と第\ref{part:catalog}部
\item \textbf{理論量が実際にどの程度の値をとるかに関心がある場合}：第\ref{ch:instantiation}章
\item \textbf{調査方法に関心がある場合}：第\ref{part:method}部と第\ref{part:results}部
\end{itemize}

各節には可能な限り数値例を付した。記号一覧は付録\ref{app:notation}にある。

\section{本稿が主張しないこと}

誤読を避けるため、先に範囲を限定する。

第一に、本稿は\textbf{予測理論ではない}。第\ref{part:results}部で事前に固定した含意は、
その多くが棄却された。本稿の枠組みは記述の道具として機能するが、
記述された量が何を予測するかについて、本稿は確立した結果を持たない。

\textbf{それにもかかわらず本稿は22の含意を事前に固定している。} 理由は二つある。

一つは、記述の枠組みが\textbf{観測されたどんな事実にも事後の説明を与えられる}ためである。
説明できることは、枠組みが正しいことの証拠にならない。
観測の前に含意を文章として固定し、外れたものを外れたまま残さなければ、
枠組みが働いた場合と、自分が後から筋を通した場合とを区別できない。
実際に後付けの物語を採用しかけた箇所は第\ref{sec:prereg-worked}節に挙げる。
\textbf{固定した含意は枠組みの主張ではなく、自分の解釈に課した拘束である。}

もう一つは、\textbf{含意が外れることと枠組みが棄却されることが別}だからである。
$\CCC$、$\kappa$、$\Phi$ という記述が機能するかどうかと、
それらの量が何を予測するかは、別の問いである。
外れた含意は後者についての推測を否定するが、前者を否定しない
（第\ref{sec:two-rejections}節）。

第二に、本稿は\textbf{規範的でない}。どの $\Phi$ が望ましいかを論じず、
どの $\Phi$ が制約の下で実現可能かのみを論じる。

第三に、\textbf{個々の構成要素の多くは既存理論に対応がある}。
付録\ref{app:priorwork}を三層に分け、
各項を説明する既存理論、それらを接続する部分、本稿の整理を区別した。
本稿の寄与は接続にあり、要素の新規性にはない。

%============================================================
\chapter{ビジネスモデル：\texorpdfstring{$\Phi$}{Φ} とキャッシュフローの生成}
\label{sec:phi}
%============================================================

\section{基本集合}

\begin{definition}[基本集合]
以下の四つを基本集合として置く。
\begin{align}
N &= \{0,1,\dots,n\} && \text{主体の集合。}\ 0\ \text{は対象企業} \\
T &= [0,H] \subset \R_{\geq 0} && \text{時間} \\
(\Omega,\mathcal{F},\{\mathcal{F}_t\}_{t\in T},\Prob) && &\text{状態空間とフィルトレーション} \\
\mathcal{D} && &\text{履行（財・サービス）の空間}
\end{align}
主体の内訳は、顧客 $C$、仕入先 $S$、資本提供者 $K$、労働 $L$、
リスク引受の第三者 $R$ を想定する。
これらは $N$ の部分集合であり、$N\setminus\{0\}$ を覆う。
個々の主体は $i\in C$ のように $N$ の元として書く。
類別をそのまま添字にしたときは、その類別に属する主体についての和を表す。
たとえば $\kappa_C=\sum_{i\in C}\kappa_i$ である。
$R$ は自由度(3) が支払主体を移す先であり、第\ref{sec:three-ops}節の移転にあたる。

規制や商慣行は現金の相手方ではないため $N$ に含めない。
これらは $\Phi$ が選べる範囲を外側から限る制約であり、
第\ref{sec:layers}節の制度層で扱う。
\end{definition}


\section{スケジュールと空間}
\label{sec:schedules}

主体 $i$ との関係を、二つのスケジュールの組として記述する。
\begin{align}
\delta_i : T\times\Omega \to \mathcal{D} && \text{履行（企業から $i$ への価値の引渡し）} \\
\pi_i : T\times\Omega \to \R && \text{決済（$i$ から企業への支払い）}
\end{align}
評価写像 $v:\mathcal{D}\to\R$ により履行を貨幣単位で測る。

これらの組の全体を空間として置く。
\begin{equation}
\Delta = \bigl\{\, (\delta_i)_{i\in N} \ \bigm| \ \delta_i : T\times\Omega\to\mathcal{D} \,\bigr\},
\qquad
\Pi = \bigl\{\, (\pi_i)_{i\in N} \ \bigm| \ \pi_i : T\times\Omega\to\R \,\bigr\}
\label{eq:spaces}
\end{equation}
$\Delta$ を\emph{履行スケジュールの空間}、$\Pi$ を\emph{決済スケジュールの空間}と呼ぶ。

\section{履行から決済への写像}

\begin{definition}[ビジネスモデル]
式\eqref{eq:spaces}の $\Delta$ と $\Pi$ を用いて、
ビジネスモデルを写像
\begin{equation}
\Phi : \Delta \longrightarrow \Pi
\label{eq:phi}
\end{equation}
として定義する。
\end{definition}

\begin{remark}[$\Phi$ と $\phi$ は別の対象である]
\label{rem:Phi-vs-phi}
大文字の $\Phi$ は式\eqref{eq:phi}の写像そのもの、すなわち\textbf{構造}を指す。
小文字の $\phi$ は、その構造のもとで一期間に生じる\textbf{余剰の大きさ}を指す量であり、
第\ref{sec:surplus}章の式\eqref{eq:surplus}で三つの源泉に分解して定める。
$\Phi$ が履行と決済の対応の形を決め、$\phi$ はその形のもとで企業に残る取り分を測る。
本稿では両者を結ぶ一般の関係式を与えない。
$\Phi$ の選択が $\phi$ の各項をどう動かすかは、
第\ref{sec:surplus}章で項ごとに個別に論じる。
\end{remark}


\section{キャッシュフローの生成}
\label{sec:cashflow}

式\eqref{eq:phi}の $\Phi$ が定まれば、主体ごとの決済が定まる。これをキャッシュフローと呼ぶ。

\begin{definition}[キャッシュフロー]
履行 $\delta\in\Delta$ を一つ定めると、$\Phi$ はその像として決済の組
$\pi = \Phi(\delta) \in \Pi$ を返す。$\Pi$ の元は主体ごとの関数の組であるから、
その第 $i$ 成分 $\pi_i$ は $T\times\Omega$ 上の関数である。これを
\begin{equation}
X(i,t,\omega) \ \equiv\ \pi_i(t,\omega), \qquad \pi = \Phi(\delta)
\label{eq:X}
\end{equation}
と書く。$X(i,t,\omega)$ は、状態 $\omega$ において時刻 $t$ に主体 $i$ との間で生じる
正味の現金移動を表す。符号は企業への流入を正とする。
仕入先や労働への支払いは $\pi_i<0$ として表される。
\end{definition}

\textbf{$\Phi$ が受け取るのは履行の組 $\delta$ ひとつであり、
$t$ や $\omega$ を引数に取らない。} 時刻と状態への依存は、
$\Phi$ が返した $\pi_i$ が持つ。

\textbf{$X$ は $\Phi$ の像である。} $\Phi$ を差し替えれば、
同じ履行 $\delta$ に対して異なる $X$ が生じる。これが本稿の中心的な操作である。

式\eqref{eq:X}は「誰が・いつ・どの状態で」の三つの添字を持つ。
通常の財務諸表は、この対象を添字方向に潰した射影である。

\begin{center}
\begin{tabularx}{\textwidth}{@{}llX@{}}
\toprule
財務諸表 & 操作 & 内容 \\
\midrule
損益計算書 & $t$ 方向の集計 & 発生主義による調整を含む \\
貸借対照表 & $t$ を固定した断面 & 未決済の残高 \\
キャッシュフロー計算書 & $t$ 方向の差分 & 現金の実際の移動 \\
\bottomrule
\end{tabularx}
\captionof{table}{財務諸表は式\eqref{eq:X}の三重添字を潰した射影である。}
\label{tab:statements-projection}
\end{center}

\begin{remark}[本稿の量はすべて水準である]
\label{rem:level-vs-change}
式\eqref{eq:X}以降に導入する量は、$\kappa_i$、$\CCC$、$\phi$ を含めてすべて水準である。
時点 $t$ における値として定義され、変化を扱う記法を持たない。

これは記述上の選択ではなく\textbf{制約である}。
第\ref{part:results}部で繰り返し確認されるとおり、
これらの量の水準は業種や商品形態によって二桁の差を持ち、
\textbf{横断比較は解釈できない}。読めるのは同一対象内の時系列変化のみである。

本稿は変化を扱う理論的な道具を持たず、
実証の各局面で個別に対処している（第\ref{sec:level-change-gap}節）。
\end{remark}

\begin{example}[三重添字の具体化]
\label{ex:triple-index}
月額1,000円のサービスを100名に提供する事業を考える。
顧客 $i\in C$ が毎月末に支払う場合、
\[
X(i,\ t,\ \omega) = \begin{cases} 1{,}000 & t \in \{1,2,\dots\}\ \text{（月末）} \\ 0 & \text{otherwise}\end{cases}
\]
であり、$\omega$ には依存しない。一方、解約が生じうるならば
$\omega$ は解約時刻を含み、$X$ は $\omega$ に依存する。

仕入先が一社であれば $S=\{s\}$ であり、サーバ費用を月額30,000円として
$X(s,t,\omega)=-30{,}000$ である。類別で書けば
$\sum_{i\in S}X(i,t,\omega)=-30{,}000$ に等しい。
\end{example}


\section{\texorpdfstring{$\Phi$}{Φ} の自由度}
\label{sec:dof}

式\eqref{eq:phi}の $\Phi$ を規定する自由度を六つに整理する。
式\eqref{eq:X}の三つの添字に対応するものと、そうでないものに分かれる。

\paragraph{添字を動かすもの。}
\begin{enumerate}[label=(\arabic*)]
\item \textbf{時間関係}：$\operatorname{supp}(\pi)$ と $\operatorname{supp}(\delta)$ の前後関係。$t$ を動かす。$\kappa$ の符号を決める。
\setcounter{enumi}{2}
\item \textbf{支払主体}：受益者本人か、第三者か、反対側の市場か。添字 $i$ を動かす。
\item \textbf{状態依存性}：$\pi$ が $\omega$ に可測か。$\omega$ を動かす。リスクの所在を決める。
これは第\ref{sec:info}章の $\mathcal{G}$ 可測性そのものであり、
$\Phi$ の設計者が自由に選べる範囲は $\mathcal{G}$ が限る
（命題\ref{prop:implementable}、注\ref{rem:G-limits-dof4}）。
\end{enumerate}

\paragraph{写像の関数形を決めるもの。}
\begin{enumerate}[label=(\arabic*)]
\setcounter{enumi}{1}
\item \textbf{依存形式}：$\pi$ が $\delta$ にどう依存するか。
線形な形式は一つの族にまとめられる。
\begin{equation}
\pi = F + p\,\delta
\label{eq:two-part}
\end{equation}
$F$ は履行量によらない部分、$p$ は単価である。
$p=0$ が定額、$F=0$ が従量、いずれも $0$ でない場合が二部料金にあたる。
線形でない形式として成果連動 $\pi = h(y)$ がある。
$y$ は成果、$h$ は成果を決済に変える関数であり、
$y$ が観測できるかは自由度(4) に属する（第\ref{sec:info}章）。
\end{enumerate}

\begin{remark}[$F$ の書体と $c$ の用法]
\label{rem:dof2-symbols}
式\eqref{eq:two-part}の $F$ は料金の固定部分であり、
状態空間の $\sigma$ 加法族 $\mathcal{F}$ とは別の量である。書体で区別する。

$c$ は本稿では限界費用を表し、つねに引数を伴って $c(\delta)$ と書く（定義\ref{def:cost}）。
定額料金には $F$ を用いる。
\end{remark}

\paragraph{定義域 \texorpdfstring{$\Delta$}{Δ} の側の性質。}
\begin{enumerate}[label=(\arabic*)]
\setcounter{enumi}{4}
\item \textbf{権利と行使の分離}：契約上の権利 $\dbar$ と実現利用 $\delta$ が乖離しうるか。
\item \textbf{反復性}：単発か、自動更新か、コミット期間つきか。取引頻度 $\nu$ を決める。
\end{enumerate}

\begin{remark}[三群の区別]
\label{rem:dof-groups}
(1)(3)(4) は式\eqref{eq:X}の添字を動かす。すなわち
\textbf{同じ $\delta$ に対して $X$ の置き場所を変える}操作である。
(2) は写像そのものの形を決め（注\ref{rem:dof2-symbols}）、(5)(6) は定義域 $\Delta$ の側の性質である。

第\ref{sec:three-ops}節のリスク操作の三分類は、(1)(3)(4) を動かす操作にあたる。
\end{remark}

このうち (1) と (2) が、$\kappa$ の符号と認知余剰の生じる余地をそれぞれ決める。
この二軸で $\Phi$ の空間を粗く分割できる（図\ref{fig:quadrant}）。
\textbf{二つは別の群から取っている。} (1) は添字を動かし、(2) は関数形を決める。

\textbf{分割の二軸と、族の判定に用いる自由度は別である。}
第\ref{sec:type-assignment}節の判定順序が用いるのは (3)(4)(6) であり、
(1) と (2) は現れない。前者は $\Phi$ の空間を粗く見るための軸、
後者は観測された $\Phi$ を一意に割り当てるための規則であって、目的が異なる。

\begin{figure}[htbp]
\centering
\insertfig{kappa-schedule}
\caption{$\Phi$ の二軸分割。縦軸が $\kappa$ の符号を、横軸が認知余剰の生じる余地を決める。}
\label{fig:quadrant}
\end{figure}

\section{リスク操作の三分類}
\label{sec:three-ops}

添字が三つであることから、$X$ に対する操作も三つに限られる。
この対応づけは本稿の整理である。

\begin{center}
\begin{tabular}{@{}clll@{}}
\toprule
操作する添字 & 名称 & 不変量 & 例 \\
\midrule
$t$ & 繰延 (timing) & 総額 & 前受金、割賦、リース \\
$i$ & 移転 (transfer) & 総額と分散 & ファクタリング、保険購入 \\
$\omega$ & 集約 (pooling) & 総額 & 保険引受、分散投資 \\
\bottomrule
\end{tabular}
\captionof{table}{リスク操作の三分類。操作する添字と不変量、代表例。}
\label{tab:three-ops}
\end{center}

この三つは、注\ref{rem:dof-groups}のとおり自由度 (1)(3)(4) を動かす操作である。
繰延が $t$ を、移転が $i$ を、集約が $\omega$ を動かす。

三つのうち、分散を実際に減少させるのは集約のみである。
これは確率論の帰結であり、$\Phi$ の構造とは独立に成立する。
\textbf{本稿が導いたものではなく、引用である。}

\begin{proposition}[集約による分散低減]
\label{prop:pooling}
分散 $\sigma^2$、相互相関 $\rho$ の同分布リスク $X_1,\dots,X_n$ に対し、
\textbf{$\rho$ が $n$ に依存しないとき}、
\begin{equation}
\operatorname{Var}\!\left[\frac{1}{n}\sum_{j=1}^{n} X_j\right]
= \frac{\sigma^2}{n}\bigl(1+(n-1)\rho\bigr)
\xrightarrow[n\to\infty]{} \rho\,\sigma^2 .
\label{eq:pooling}
\end{equation}
\end{proposition}

\begin{proof}
共分散の双線形性より
\[
\operatorname{Var}\Bigl[\sum_j X_j\Bigr]
= \sum_j \operatorname{Var}[X_j] + \sum_{j\neq k}\operatorname{Cov}[X_j,X_k]
= n\sigma^2 + n(n-1)\rho\sigma^2 .
\]
両辺を $n^2$ で割れば第一の等式を得る。
$n\to\infty$ のとき $\frac{\sigma^2}{n}\to 0$、
$\frac{(n-1)\rho\sigma^2}{n}\to\rho\sigma^2$ であるから極限が従う。
\end{proof}

\begin{remark}[独立性は有限の資源である]
\label{rem:independence-finite}
式\eqref{eq:pooling}の極限は $\rho\sigma^2$ であり、$\rho>0$ である限りゼロにならない。
集約は独立性を消費する操作であり、独立性は有限の資源である。
地震保険の引受制限、プラットフォームの景気感応性、
金融機関のシステミックリスクは、いずれもこの式の同じ項に由来する。
\end{remark}

\begin{remark}[$\rho$ の内生性]
\label{rem:rho-endogenous}
命題\ref{prop:pooling}は $\rho$ を外生的な定数とした。この仮定は成立しにくい。

$\rho$ が内生的になる機構は、集約の対象が枯渇することではなく、
\textbf{集約という行為そのものが相関を生む}ことにある。
\cite{wagner2010} の例が明快である。
二つの主体が二つのリスクを等分に引き受けて完全に分散すると、
各主体の個別の破綻確率は下がるが、
\textbf{両者のポートフォリオは完全に相関する}。
リスクの性質は変わらないのに、主体間の相関が生じる。

したがって $\rho$ は集約者の数と保有の重なりの関数であり、
式\eqref{eq:pooling}の極限 $\rho\sigma^2$ は集約が進むほど大きくなりうる。
証明は形式的には正しいが、
\textbf{適用範囲は集約者が少数で保有の重なりが小さい範囲に限られる}。
本稿はこの依存を扱わない（第\ref{sec:independence-assumption}節）。
\end{remark}

\begin{remark}[集約の社会的最適性]
\label{rem:pooling-social}
注\ref{rem:rho-endogenous}の機構は、個別と全体で最適が食い違うことを含意する。
\cite{ibragimov2011} は、仲介者にとって最適な分散が社会的には最適でなく、
リスク分担を制限したほうが望ましい条件を導いている。

本稿は個別事業者の視点で書かれており、
$\Phi$ の選択が系に与える影響を扱わない。
第\ref{part:catalog}部の族5-4（保険引受）と族6-2（エスクロー）は
集約を業とする類型であり、この論点が直接効く領域である。
\end{remark}

\begin{example}[相関が効く速さ]
$\sigma=100$、$\rho=0$ のとき、$n=100$ で標準偏差は $100/\sqrt{100}=10$ に下がる。
一方 $\rho=0.3$ では
\[
\operatorname{Var}=\frac{100^2}{100}\bigl(1+99\times 0.3\bigr)=305.5,\qquad
\text{標準偏差}\approx 17.5 .
\]
$n\to\infty$ でも標準偏差は $\sqrt{0.3}\times 100 \approx 54.8$ より下がらない。
相関 $0.3$ という穏当な値でも、分散低減の効果は半分以下に落ちる。
\end{example}

\subsection{技術が定める費用}
\label{sec:production}

履行 $\delta$ を生むには費用がかかる。本稿は投入の内訳を扱わないので、
生産の技術は\textbf{履行から費用への対応 $C_{\mathrm{prod}}$ としてのみ現れる}。

\textbf{$C_{\mathrm{prod}}$ は $\Phi$ とは独立に与えられる。}
同じ技術に対して異なる $\Phi$ を置くことができ、
逆に同じ $\Phi$ を異なる技術の上に置くこともできる。
本稿が動かすのは $\Phi$ であって技術ではない。

\begin{definition}[費用の分割]
\label{def:cost}
総費用を次の二つに分割する。
\begin{equation}
C(\delta,\Phi) = \underbrace{C_{\mathrm{prod}}(\delta)}_{\text{技術が定める}}
 + \underbrace{C_{\mathrm{admin}}(\Phi,\delta)}_{\text{契約が生む}}
\label{eq:cost-split}
\end{equation}
第一項について
\begin{align}
C_{\mathrm{prod}}(0) && &\text{固定費} \\
c(\delta) = \frac{dC_{\mathrm{prod}}}{d\delta} && &\text{限界費用}
\end{align}
と書く。
\end{definition}

\begin{remark}[$c(\delta)$ に形状を仮定しない]
本稿は $c(\delta)$ の単調性も凸性も仮定しない。
費用構造は業種ごとに大きく異なり、一つの形状を置くと適用範囲が狭まるためである。

代償として、\textbf{最適な履行量 $\delta$ を内点で定めることができない}。
式\eqref{eq:cap-simplified}の容量制約は上限を与えるが、
その手前で費用と価格が釣り合う点は本稿の仮定からは出ない。
以降で $\delta$ の水準が問題になる箇所（第\ref{ch:solo}章の配分など）では、
制約の端点として扱う。
\end{remark}

この分割に意味があるのは、第二項が実在するからである。
$C_{\mathrm{admin}}$ は $\Phi$ の六自由度（第\ref{sec:dof}節）のすべてから生じる。

\begin{center}
\begin{tabularx}{\textwidth}{@{}llX@{}}
\toprule
種別 & 経路 & 内容 \\
\midrule
設定 & 締結 & $\Phi$ を相手と合意する費用。定型契約では小さく、個別交渉では大きい \\
運用 & 測定 & $\pi=p\,\delta$ では $\delta$ の計測と課金が要る \\
 & 権利管理 & $\dbar$ の管理。会員資格、残高の追跡 \\
 & 継続処理 & 自動更新、解約処理 \\
履行確保 & 回収・貸倒れ & $\tau_C>0$ では請求と督促、回収不能の損失 \\
 & 検証 & $\pi=h(y)$ では $y$ の確認 \\
 & 紛争 & 契約に書けても争いは生じる \\
外部委託 & 仲介 & 決済、配信、集客の代行 \\
\bottomrule
\end{tabularx}
\captionof{table}{$C_{\mathrm{admin}}$ が生じる経路。}
\label{tab:cadmin-routes}
\end{center}

\begin{remark}[前受にも費用がかかる]
$\tau_C<0$（前受）は回収費用を要しないが、
預り金の管理、返金対応、規制対応が生じる。
第\ref{sec:bootstrap}節で扱う資金決済法の供託義務がその例である。
\textbf{$\tau$ の符号によらず $C_{\mathrm{admin}}>0$ である。}
\end{remark}

\begin{remark}[$C_{\mathrm{admin}}$ を無視する条件]
\label{rem:cadmin-negligible}
本稿は以下、$C_{\mathrm{admin}}$ を明示しない。無視できる条件は二つある。

\begin{enumerate}[label=(\arabic*)]
\item \textbf{運用がシステム化されている。}
測定・権利管理・継続処理は、仕組みが構築されていれば
取引あたりの限界費用がゼロに近づく。
\item \textbf{仲介を用いない。}
仲介費用は取引額に比例し、システム化では消えない。
第\ref{sec:platform-fee}章はこの費用を測定し、
注\ref{rem:ladder-is-cadmin}で $C_{\mathrm{admin}}$ の種別に対応づける。
\end{enumerate}

条件が成立する範囲は限定的である。
第\ref{ch:solo}章の一人事業と第\ref{part:results}部の族2では成立するが、
\textbf{第\ref{part:industry}章が扱う小規模企業では成立しない}。
請求と督促が人手で行われる場合、$C_{\mathrm{admin}}$ は無視できない。

この省略の代償を明示しておく。$C_{\mathrm{admin}}$ を落とすため、
本稿は\textbf{$\Phi$ の選択が費用を通じて $\phi$ に跳ね返る経路を扱えない}。
式\eqref{eq:cost-split}は以降で使わず、$C_{\mathrm{prod}}$ のみが現れる。
$\tau_i$ を延ばす契約変更が回収費用を増やす、といった対応関係は、
本稿の枠組みの中では書けない。
\end{remark}

以上より、費用構造が同一でも異なる契約形態を選べるという自由度が、
ここでいうビジネスモデルにあたる。
ただしこの自由度は完全ではなく、$\Phi$ の選択は $C_{\mathrm{admin}}$ を通じて費用に跳ね返る。
$C_{\mathrm{prod}}$ が現れる箇所は注\ref{rem:cost-usage}にまとめた。

\begin{remark}[$\Phi$ の切り替え費用と待機の価値]
\label{rem:switching-cost}
本稿は $\Phi$ を選択変数として扱うが、
これは\textbf{切り替え費用をゼロと置くこと}に等しい。
一度選んだ $\Phi$ を別の $\Phi$ に変える費用は $C_{\mathrm{admin}}$ の設定項に属し、
一般にゼロではない。

不可逆投資の理論（\cite{dixit1994}）によれば、
決定が不可逆であるのは物理的に取り消せないためではなく、
\textbf{取り消しの費用が高く特定の経路に固定されるため}である。
切り替え費用を $C_{\mathrm{sw}}$、$\Phi$ の変更で得られる余剰の増分を $\Delta\phi$ とすると、

\begin{center}
\begin{tabularx}{\textwidth}{@{}lX@{}}
\toprule
$C_{\mathrm{sw}} = \infty$ & 変更できない。最初の選択が確定的 \\
$C_{\mathrm{sw}} \ll \Delta\phi$ & 本稿の仮定。$\Phi$ は自由に選べる \\
$C_{\mathrm{sw}} \sim \Delta\phi$ & 選択を遅らせることに価値がある \\
\bottomrule
\end{tabularx}
\captionof{table}{切り替え費用 $C_{\mathrm{sw}}$ の大小と $\Phi$ 選択の自由度。}
\label{tab:switching-cost-cases}
\end{center}

第三の場合には\textbf{待機の価値}が生じる。
不確実性が解消するまで $\Phi$ を確定させないことが最適になりうる。
本稿はこの含意を扱わない。第\ref{ch:solo}章の一人事業では
切り替え費用が相対的に大きいため、この論点は無視できない可能性がある。
\end{remark}

\begin{remark}[$C$ が現れる箇所]
\label{rem:cost-usage}
定義\ref{def:cost}の量は以降で繰り返し用いる。
$\phi_{\mathrm{prod}}$ は $c(\delta)$ の優位として定義し（第\ref{sec:surplus}章）、
限界費用がゼロに近いことが容量制約を緩める条件となり（第\ref{ch:solo}章）、
$C$ と $\Phi$ がともに事前指定できることが複製可能性の条件となる（第\ref{part:unmanned}章）。
物理層はこの $C$ の性質を指す（第\ref{sec:layers}節）。
\end{remark}

\begin{example}[同一の $C_{\mathrm{prod}}$、異なる $\Phi$]
同一の業務用ソフトウェアを、次の三通りで提供できる。
\begin{center}
\begin{tabular}{@{}llr@{}}
\toprule
$\Phi$ & 決済 & $\kappa_C$（期央、顧客1名） \\
\midrule
買切り & 導入時に240,000円 & \textminus 120,000 円 \\
年額 & 毎年初に60,000円 & \textminus 30,000 円 \\
月額後払 & 毎月末に5,000円 & +2,500 円 \\
\bottomrule
\end{tabular}
\captionof{table}{同一の $C_{\mathrm{prod}}$ に対する三通りの $\Phi$ と $\kappa_C$。}
\label{tab:same-f-different-phi}
\end{center}
技術は同一であり、限界費用も同一である。異なるのは $\Phi$ のみだが、
$\kappa$ の符号が反転し、必要な運転資本が変わる。
\end{example}

\section{権利と行使の乖離}
\label{sec:breakage}

自由度 (2) と (5) は連動している。$\pi$ が $\delta$ に比例する場合、
顧客は使った分だけ支払うため、権利と行使は一致する。
$\pi$ が $\delta$ と独立な場合にのみ、両者が乖離しうる。

式\eqref{eq:phi}の自由度(5) が立つとき、次を定める。

\begin{definition}[権利と行使]
\label{def:breakage-def}
契約が顧客に与える上限を $\dbar$、実現する利用を $\delta$ とする。
乖離量を $\dbar-\delta \geq 0$ と書き、これに由来する粗利を
\begin{equation}
\phi_{\mathrm{cog}} = p\cdot\E[\dbar-\delta]
\label{eq:cog}
\end{equation}
と定義する。
\end{definition}

\begin{remark}[乖離が生じる契約形態]
\label{prop:breakage}
式\eqref{eq:cog}が正となりうるのは、契約が $\dbar$ を明示的に定める場合に限られる。

$\pi(t)=p\,\delta(t)$ の従量契約では、顧客の支払総額は $p\int\delta$ である。
利用しなければ支払いも生じないため、契約は上限 $\dbar$ を定める必要がなく、
差 $\dbar-\delta$ が定義されない。

$\pi$ が $\delta$ に依存しない定額契約では、
対価が利用量と無関係に定まるため、契約は権利の範囲 $\dbar$ を明示する。
このとき $\delta<\dbar$ が実現しうる。
二部料金 $\pi=F+p\,\delta$ では、固定部分 $F$ に対応する範囲についてのみ乖離が生じる。

\textbf{これは定義\ref{def:breakage-def}の帰結であり、独立した主張ではない。}
$\dbar$ が外生に与えられることを前提とする。
顧客がメニューから $\dbar$ を選ぶ場合は同時決定となり、
差を取ることの意味が変わる（第\ref{sec:independence-assumption}節）。
\end{remark}

\begin{remark}[容量制約の一般形と本稿の簡略化]
\label{rem:cap-simplification}
供給側の容量には、施設の収容人数、設備の台数、人員の可処分時間などが該当する。
第\ref{sec:phi}章の定義より $\delta_i$ は $T\times\Omega$ 上の関数であるから、
容量制約の一般形は各時刻で評価される。
\begin{equation}
\sum_i \delta_i(t,\omega) \ \leq\ \mathrm{Cap}(t,\omega)
\qquad \forall t\in T,\ \forall\omega\in\Omega
\label{eq:cap-general}
\end{equation}
容量それ自体が時刻に依存するのは、営業時間、設備の稼働時間帯、
人員の配置が時間によって変わるためである。

\textbf{本稿では以下、$\mathrm{Cap}$ を定数として扱う。}
すなわち式\eqref{eq:cap-general}を期間全体で積分した
\begin{equation}
\sum_i \delta_i \ \leq\ \mathrm{Cap}
\label{eq:cap-simplified}
\end{equation}
の形で用いる。$\delta_j$ は総量を表す。

この簡略化は、需要の時間分布が会員間で概ね等しく、
かつ容量が時間帯によらず一定である場合に正当化される。
両者が成立しない場合、拘束するのはピーク時刻のみであり、
式\eqref{eq:cap-simplified}は必要以上に強い条件となる。
第\ref{part:industry}章では、この簡略化が現実と合わない観測を扱う。
\end{remark}

\begin{proposition}[容量制約下の会員制]
\label{prop:capacity}
式\eqref{eq:cap-simplified}の容量制約が存在し、
\textbf{会員数 $n$ と平均利用 $\E[\delta]$ が独立に決まるとき}（$m$ は粗利率に用いるため会員数は $n$ と書く）、
定額会員制 $\pi_j=F$ が実行可能であるためには
\begin{equation}
\sum_{j}\E[\delta_j] \ \ll\ \sum_{j}\dbar_j
\end{equation}
が必要である。すなわち低稼働が実行可能性の条件になる。
\end{proposition}

\begin{proof}
会員数を $n$、各会員の契約上の上限を $\dbar$ とする。
全会員が権利を行使すれば需要は $n\dbar$ となる。
実行可能性は $n\,\E[\delta] \leq \mathrm{Cap}$ を要求する。
一方、会員が $\dbar$ に価値を感じて加入している以上 $n\dbar > \mathrm{Cap}$ が通常成立する
（そうでなければ容量制約が拘束していない）。
両者を合わせると $\E[\delta] \leq \mathrm{Cap}/n < \dbar$ を得る。
\end{proof}

\begin{example}[二種の会員制]
定員100名の施設が1,000名の会員を持ち、月会費8,000円を課すとする。
$\dbar$ を「毎日利用可能」とすれば $n\dbar$ は容量を大きく超える。
実際の平均利用率が10\%であれば $n\E[\delta]=100$ で制約はぎりぎり満たされる。
\textbf{利用率が30\%に上昇すると事業は物理的に破綻する。}

対照的に、会費が仕入割引の権利を与える形態では容量制約が存在せず、
利用率の上昇は仕入規模の拡大を通じて有利に働く。
外形上は同じ族2-1 に属するが、実行可能領域の形が異なる。
\end{example}

%============================================================
\chapter{信用ポジション：\texorpdfstring{$\kappa$}{κ} と遅れ時間 \texorpdfstring{$\tau$}{τ}}
\label{sec:kappa}
%============================================================

\section{履行と決済の差}

第\ref{sec:schedules}節の履行 $\delta_i$ と、
式\eqref{eq:phi}の $\Phi$ が定める決済 $\pi_i$（ただし $\pi=\Phi(\delta)$）の累積を
\begin{equation}
D_i(t) = \int_0^t v(\delta_i(s))\,ds, \qquad
P_i(t) = \int_0^t \pi_i(s)\,ds
\label{eq:cum}
\end{equation}
と書く。$D_i$ は与えた価値、$P_i$ は受け取った対価である。

\begin{definition}[信用ポジション]
主体 $i$ に対する信用ポジションを
\begin{equation}
\kappa_i(t) = D_i(t) - P_i(t)
\label{eq:kappa}
\end{equation}
と定義する。式\eqref{eq:cum}より
\[
\kappa_i(t) = \int_0^t v(\delta_i(s))\,ds - \int_0^t \pi_i(s)\,ds
\]
であり、\textbf{$\kappa_i$ は式\eqref{eq:phi}の $\Phi$ の関数である。}
同じ履行 $\delta$ に対して $\Phi$ を差し替えれば $\kappa_i$ が変わる。$\kappa_i>0$ は企業が $i$ に与信している状態（履行済み・未回収）、
$\kappa_i<0$ は $i$ が企業に与信している状態（受領済み・未履行）を表す。
\end{definition}

信用ポジションの総和に棚卸資産を加えた量を、以降で繰り返し用いる。

\begin{definition}[運転資本]
\label{def:wc}
棚卸資産を $\mathrm{Inv}(t)\geq 0$ として
\begin{equation}
W(t) = \sum_{i\in N}\kappa_i(t) + \mathrm{Inv}(t)
\label{eq:wc}
\end{equation}
と定める。$\kappa_i$ が $\Phi$ の関数であるから、$W$ も $\Phi$ の関数である。
\end{definition}

式\eqref{eq:wc}の $W$ は事業に拘束されている資金の量であり、
第\ref{sec:growth}章ではこれを売上速度で割って $\CCC$ を定義する。

式\eqref{eq:kappa}の利点は統合性にある。

\begin{center}
\begin{tabular}{@{}llll@{}}
\toprule
勘定科目 & 主体 & 符号 & 意味 \\
\midrule
売掛金 & $C$ & $\kappa_C>0$ & 企業が顧客に与信 \\
前受金・契約負債 & $C$ & $\kappa_C<0$ & 顧客が企業に与信 \\
買掛金 & $S$ & $\kappa_S<0$ & 仕入先が企業に与信 \\
前渡金 & $S$ & $\kappa_S>0$ & 企業が仕入先に与信 \\
\bottomrule
\end{tabular}
\captionof{table}{式\eqref{eq:kappa}が統合する四つの勘定科目。}
\label{tab:kappa-accounts}
\end{center}

四つの別勘定ではなく、符号と添字の異なる単一の量である。

\subsection{各項を説明する既存理論}
\label{sec:kappa-literature}

$\kappa_i$ の各項には、それぞれ確立した理論がある。
本稿はこれらを土台として用いる。

\begin{center}
\begin{tabularx}{\textwidth}{@{}llX@{}}
\toprule
主体と符号 & 現象 & 説明する理論 \\
\midrule
$\kappa_C>0$（企業間） & 掛売 & 取引信用の四類型。金融優位、価格差別、品質保証、清算優位 \cite{petersen1997,schwartz1974} \\
$\kappa_C<0$（企業間） & 前受 & reverse trade credit。供給者金融、取引保証、交渉力 \cite{mateut2014,daripa2011} \\
$\kappa_C<0$（対消費者） & 前払い契約 & メンタルアカウンティング、退蔵の実証 \cite{nber2026pnbl} \\
$\kappa_S<0$ & 買掛 & 取引信用の鏡像 \\
$\kappa_L<0$ & 給与後払い & bonding argument、後払いによる規律 \cite{akerlof1986} \\
$\kappa_K<0$ & 借入 & 企業金融の標準理論 \\
$\kappa_K^{i}$ & 出資（残余請求権） & 主体依存。定義\ref{def:kappa-K} \\
\bottomrule
\end{tabularx}
\captionof{table}{$\kappa_i$ の各項を説明する既存理論。}
\label{tab:kappa-literature}
\end{center}

\begin{remark}[四つの理論は互いを参照しない]
\label{rem:disjoint-literatures}
表\ref{tab:kappa-literature}の理論は、それぞれ異なる分野で異なる問いに答えるために発展した。

\begin{center}
\begin{tabularx}{\textwidth}{@{}llX@{}}
\toprule
理論 & 分野 & 中心的な問い \\
\midrule
取引信用 & 企業金融 & なぜ供給者が銀行より安く与信できるか \\
bonding argument & 労働経済学 & なぜ非自発的失業が存在するか \\
シグナリング & 情報の経済学 & 能力をどう開示するか \\
前払い契約の研究 & 消費者行動 & なぜ消費者が先に払うか \\
\bottomrule
\end{tabularx}
\captionof{table}{$\kappa_i$ の各項を説明する理論の分野と中心的な問い。}
\label{tab:disjoint-theories-questions}
\end{center}

問いが異なるため、これらは互いを引用しない。
しかし式\eqref{eq:cash}において、
\textbf{顧客の前払い、仕入先の掛売、労働者の後払い、資本の出資は
すべて同一の和の項として現れる}。

各理論は和の一項ずつを説明するが、和として扱う枠組みを持たない。
本稿の寄与はこの統合にあり、個々の項の説明にはない。
\end{remark}

累積の関係を図\ref{fig:kappa}に示す。

\begin{figure}[htbp]
\centering
\insertfig{phi-quadrant}
\caption{累積履行 $D(t)$ を基準線とした信用ポジション。決済曲線 $P(t)$ が $D(t)$ の上にあれば $\kappa<0$、下にあれば $\kappa>0$。}
\label{fig:kappa}
\end{figure}

\subsection{取引の両側：$\kappa$ は反対称である}
\label{sec:kappa-antisym}

式\eqref{eq:kappa}は自社の帳簿で書いた量である。
同じ取引を相手方 $i$ の側で書くと、与えた価値と受け取った対価が入れ替わる。

\begin{proposition}[反対称性]
\label{prop:kappa-antisym}
二つの主体 $i$ と $j$ が評価写像 $v$ を共有するとき、
$i$ の帳簿における $j$ への信用ポジションを $\kappa_j^{(i)}$ と書けば
\begin{equation}
\kappa_j^{(i)} = -\,\kappa_i^{(j)}
\label{eq:kappa-antisym}
\end{equation}
が成立する。
\end{proposition}

\begin{proof}
$i$ が $j$ に与えた価値は $j$ が $i$ から受け取った価値であり、
$i$ が $j$ から受け取った対価は $j$ が $i$ に与えた対価である。
したがって $D_j^{(i)}=P_i^{(j)}$、$P_j^{(i)}=D_i^{(j)}$ であり、
式\eqref{eq:kappa}に代入すれば従う。
\end{proof}

\begin{remark}[$v$ の共有が要る]
式\eqref{eq:kappa-antisym}は両者が同じ $v$ を用いることを前提とする。
注\ref{rem:subjectivity}のとおり $v$ は主観的でありうるが、
決済額 $\pi$ は $\mathcal{G}$ 可測で一致するから、
不一致が生じるのは $D$ の側だけである。
検収基準の相違や係争中の債権がこれにあたる。
\end{remark}

反対称性から、集計についての帰結が出る。

\begin{corollary}[企業間の信用は集計で消える]
\label{cor:kappa-nets-out}
法人部門に属する主体の集合を $\mathcal{N}$ とすると、
$\mathcal{N}$ の内部で閉じた信用ポジションの総和はゼロである。したがって
\begin{equation}
\sum_{i\in\mathcal{N}} W^{(i)}
= \kappa_{\text{対非法人}} + \sum_{i\in\mathcal{N}}\mathrm{Inv}^{(i)}
\label{eq:aggregate-W}
\end{equation}
が成立する。ここで $\kappa_{\text{対非法人}}$ は法人部門と家計・政府・海外との
あいだの正味ポジションである。
\end{corollary}

\begin{proof}
$\mathcal{N}$ の内部の各対 $(i,j)$ について式\eqref{eq:kappa-antisym}より
$\kappa_j^{(i)}+\kappa_i^{(j)}=0$ である。
定義\ref{def:wc}の $W$ を $\mathcal{N}$ の全主体について足すと内部の対は消え、
外部との対と在庫が残る。
\end{proof}

\begin{corollary}[すべての階層が与信の出し手にはなれない]
\label{cor:no-uniform-direction}
法人部門を互いに素な階層に分割したとき、
すべての階層が企業間で正味の与信の出し手であることはありえない。
ある階層が正味の出し手であれば、別の階層は正味の受け手である。
\end{corollary}

\begin{proof}
系\ref{cor:kappa-nets-out}より、階層をまたぐ企業間ポジションの総和はゼロである。
すべての階層で正味ポジションが狭義に正であれば総和も狭義に正となり、ゼロに矛盾する。
\end{proof}

系\ref{cor:no-uniform-direction}は、規模階層の比較に構造的な制約を課す。
第\ref{part:results}部で与信の向きを規模別に測るとき、
観測される符号の一部はこの制約から自動的に決まる。

\begin{remark}[集計の $\CCC$ が測っているもの]
式\eqref{eq:aggregate-W}を売上の総和で割ると、集計の $\CCC$ が得られる。
分子から企業間ポジションが落ちているから、
\textbf{集計の $\CCC$ は企業間信用の指標ではない}。
残るのは対非法人ポジションと在庫である。
第\ref{part:industry}章で全産業の $\CCC$ を報告するが、
その水準の解釈にはこの区別が要る。
\end{remark}

\subsection{$\Phi$ の並置に対する加法性}
\label{sec:kappa-additive}

一つの主体が複数の $\Phi$ を並行して運用することがある。

\begin{proposition}[$\kappa$ の加法性]
\label{prop:kappa-additive}
互いに素な履行 $\delta^1,\dots,\delta^K$ に対して
$\Phi^1,\dots,\Phi^K$ を並行して運用するとき、
\begin{equation}
\kappa = \sum_{k=1}^{K}\kappa^{k},
\qquad
W = \sum_{k=1}^{K} W^{k}
\label{eq:kappa-additive}
\end{equation}
が成立する。
\end{proposition}

\begin{proof}
式\eqref{eq:cum}の $D_i$ と $P_i$ は積分であり、被積分関数について加法的である。
式\eqref{eq:kappa}は両者の差であるから加法性を保つ。
\end{proof}

加法的なのは $\kappa$ と $W$ であって、$\CCC$ ではない。
$\CCC=W/r$ は比であり、比は加法的でない（系\ref{cor:ccc-weighted}）。

\begin{remark}[分解は一意でない]
\label{rem:decomposition-not-unique}
式\eqref{eq:kappa-additive}は分解の存在を述べるが、一意性を述べない。
同じ $\kappa$ を与える $\{\Phi^k\}$ の組は複数ある。
第\ref{ch:catalog}章の類型は $\Phi$ の生成系であって基底ではない。
\end{remark}

\subsection{遅れ時間としての表現}
\label{sec:tau}

$\kappa_i$ は累積の差であり、その内実は\textbf{履行と決済の時間差}である。
これを明示する。

\begin{definition}[決済の遅れ]
主体 $i$ について、履行から決済までの平均的な遅れを $\tau_i$ と書く。
符号は $\kappa_i$ に合わせ、決済が履行より遅い場合を正とする。
\begin{center}
\begin{tabularx}{\textwidth}{@{}lXl@{}}
\toprule
符号 & 状態 & 例 \\
\midrule
$\tau_i>0$ & 決済が履行より遅い（企業が与信） & 売掛金、掛売 \\
$\tau_i<0$ & 決済が履行より早い（$i$ が与信） & 前受金、買掛金 \\
\bottomrule
\end{tabularx}
\captionof{table}{決済の遅れ $\tau_i$ の符号と対応する状態。}
\label{tab:tau-sign}
\end{center}
\end{definition}

式\eqref{eq:kappa}の $\kappa_i$ を、遅れ時間 $\tau_i$ で表す。

\begin{proposition}[信用ポジションの分解]
\label{prop:kappa-tau}
定常状態において履行速度が一定であるとき、時間平均として
\begin{equation}
\bar{\kappa}_i = r\,\tau_i
\label{eq:kappa-tau}
\end{equation}
が成立する。
\end{proposition}

\begin{proof}
決済が履行から一律に $\tau_i$ 遅れるとき、時刻 $t$ における未決済の累積は
直近 $\tau_i$ 期間の履行にあたる。
\[
\kappa_i(t) = \int_{t-\tau_i}^{t} r(s)\,ds .
\]
$r$ が一定であれば右辺は $r\tau_i$ に等しい。
遅れが分布を持つ場合は $\tau_i$ をその平均とすればよい。
\end{proof}

\begin{remark}[平均としてのみ成立する]
\label{rem:kappa-tau-average}
実際の決済は離散的である。月末締め翌月末払いであれば
$\pi$ は月に一度のパルスとなり、$\kappa_i(t)$ は鋸歯状に変動する。
式\eqref{eq:kappa-tau}はその時間平均について成立するにすぎない。

期末残高を用いる統計では、決算期が特定の月に集中していれば
鋸歯の位相が揃い、平均から系統的に乖離しうる。
本稿はこの偏りを補正していない。
\end{remark}

\begin{example}[年払いサブスクリプションの $\kappa$]
月額1,000円のサービスを年払い12,000円で1月1日に受け取る。
履行は12か月かけて均等に行われるから、$t$ か月経過時点で
\[
\kappa_C(t) = 1{,}000\,t - 12{,}000 .
\]
$t=0$ で $-12{,}000$、$t=6$ で $-6{,}000$、$t=12$ で $0$ である。
顧客100名なら、期央の平均で $\kappa_C = -600{,}000$ 円。
これがそのまま無利子の調達額となる。
\end{example}

\section{現金残高の分解}

式\eqref{eq:kappa}の $\kappa_i$ を用いて現金残高を分解する。

\begin{proposition}[現金の信用分解]
\label{prop:cash}
自己資本 $E(t)$、棚卸資産 $\mathrm{Inv}(t)$、固定資産 $A(t)$ を用いて
\begin{equation}
M(t) = E(t) + \sum_{i\in N}\max\{-\kappa_i(t),0\} - \sum_{i\in N}\max\{\kappa_i(t),0\}
- \mathrm{Inv}(t) - A(t)
\label{eq:cash}
\end{equation}
が成立する。
\end{proposition}

\begin{proof}
貸借対照表恒等式より、資産の総和は負債と純資産の総和に等しい。
資産を現金 $M$、正の信用ポジション（売掛金・前渡金など）、棚卸資産 $\mathrm{Inv}$、
固定資産 $A$ に分割し、負債を負の信用ポジション（買掛金・前受金など）に、
純資産を $E$ に対応させると
\[
M + \sum_i \max\{\kappa_i,0\} + \mathrm{Inv} + A
= \sum_i \max\{-\kappa_i,0\} + E .
\]
$M$ について解けば式\eqref{eq:cash}を得る。
\end{proof}

式\eqref{eq:cash}は恒等式の並べ替えにすぎないが、読み方を規定する。
第二項が他者から引き出した与信の総額、第三項が他者に与えた与信の総額である。

\begin{corollary}[現金の意味]
企業の現金とは、全ての相手方から引き出した信用の合計から、
与えた信用と固定化した資産を差し引いた残余である。
\end{corollary}

\subsection{正味の形}
\label{sec:cash-net}

式\eqref{eq:cash}の第二項と第三項は、総額を分けて見せるために置いている。
正味だけを見るなら、任意の実数 $x$ について $\max\{-x,0\}-\max\{x,0\}=-x$ であるから、
二項は打ち消し合う。

\begin{proposition}[現金は運転資本の裏返しである]
\label{prop:cash-net}
定義\ref{def:wc}の運転資本 $W$ を用いると、式\eqref{eq:cash}は
\begin{equation}
M(t) = E(t) - W(t) - A(t)
\label{eq:cash-net}
\end{equation}
と書ける。
\end{proposition}

\begin{proof}
$\max\{-\kappa_i,0\}-\max\{\kappa_i,0\}=-\kappa_i$ を $i$ について足すと
$\sum_i\max\{-\kappa_i,0\}-\sum_i\max\{\kappa_i,0\}=-\sum_i\kappa_i$ である。
式\eqref{eq:cash}に代入し、$\mathrm{Inv}$ をまとめて $W$ とすればよい。
\end{proof}

例\ref{ex:three-sources}で確かめると、$\sum_i\kappa_i=-600-50+200=-450$、
$\mathrm{Inv}=0$ より $W=-450$ であり、$100+450-30=520$ で一致する。

式\eqref{eq:cash-net}は式\eqref{eq:cash}と同じ内容だが、
\textbf{$W$ を通じて第\ref{sec:growth}章の $\CCC$ に接続する}。
第\ref{sec:objective}節の制約 $M\geq 0$ が
$\Phi$ の選択にどう効くかは、この形を経由して定まる。

\begin{remark}[既存の恒等式である]
式\eqref{eq:cash-net}は財務分析で用いられる恒等式
（現金 ＝ 自己資本 $-$ 正味営業資産）と同一である。
本稿における意味は、$W$ が式\eqref{eq:kappa}を通じて $\Phi$ の関数であり、
したがって現金制約が $\Phi$ についての制約になる点にある。
\end{remark}

\begin{example}[三者からの調達]
\label{ex:three-sources}
自己資本100万円で開業した事業が、顧客から年払いで600万円を受け取り
（$\kappa_C=-600$万円）、クラウド事業者に対して未払い50万円を持ち
（$\kappa_S=-50$万円）、大口顧客への売掛金が200万円ある（$\kappa_{C'}=+200$万円）。
在庫はなく、固定資産は30万円とする。
\[
M = 100 + (600+50) - 200 - 0 - 30 = 520\ \text{万円}.
\]
自己資本100万円に対し手元現金は520万円である。
差額420万円は顧客と仕入先から引き出した与信であり、返済期限を持つ。
\end{example}

\subsection{自己資本と資本提供者}
\label{sec:equity}

命題\ref{prop:cash}の $E$ は残余として現れるが、内訳を明示する。
式\eqref{eq:cash}はこの $E$ を用いている。

\begin{definition}[自己資本]
\label{def:equity}
払込資本を $P_K(t)$、配当の累積を $\mathrm{Div}(t)$ として
\begin{equation}
E(t) = P_K(t) + \int_0^t \phi(s)\,ds - \mathrm{Div}(t)
\label{eq:equity}
\end{equation}
と定める。第二項が内部留保である。
\end{definition}

定義\ref{def:equity}は第\ref{sec:surplus}章の $\phi$ と本章の $E$ を結ぶ。
$\phi$ の三分解を代入すれば
\[
\text{内部留保} = \int (\phi_{\mathrm{prod}} + \phi_{\mathrm{barg}} + \phi_{\mathrm{cog}})\,ds - \mathrm{Div}
\]
となる。第\ref{sec:surplus}章では $\phi_{\mathrm{cog}}$ が学習と規制で削られることを扱うが、
\textbf{削られるのは将来のフローであり、過去に蓄積した分は $E$ に残る}。

\begin{definition}[資本提供者との信用ポジション]
\label{def:kappa-K}
資本提供者 $i\in K$ の確率測度 $\Prob_i$ の下での残余請求権の価値を
$V_i(t) = \E_i[\text{将来キャッシュフローの現在価値}]$ として
\begin{equation}
\kappa_K^{i}(t) = V_i(t) - P_K(t)
\label{eq:kappa-K}
\end{equation}
と定める。
\end{definition}

\begin{proposition}[$\kappa_K$ は式\eqref{eq:cash}に入らない]
\label{prop:kappa-K-excluded}
$\kappa_K^{i}$ は資本提供者 $i$ に依存するため、式\eqref{eq:cash}の恒等式に含められない。
\end{proposition}

\begin{proof}
式\eqref{eq:cash}は貸借対照表恒等式の変形であり、
全主体が同一の値を認識する量のみで構成される。
売掛金や買掛金は当事者間で額が一致するが、
式\eqref{eq:kappa-K}の $V_i$ は $\Prob_i$ に依存し一致しない。
主体依存の量を含めれば恒等式が主体ごとに変わり、恒等式でなくなる。
\end{proof}

\begin{remark}[主観性の三段階]
\label{rem:subjectivity}
本稿には観測できない量が複数現れるが、性質が異なる。
第\ref{sec:info}章の $\mathcal{G}$ を用いて三段階に分けられる。

\begin{center}
\begin{tabularx}{\textwidth}{@{}llX@{}}
\toprule
段階 & 条件 & 例 \\
\midrule
共有 & $\mathcal{G}$ 可測 & 決済額 $\pi$、売掛金 \\
私的情報 & $\mathcal{F}_t$ 可測だが $\mathcal{G}$ 可測でない & 評価写像 $v$、努力 $a$ \\
測度依存 & $\mathcal{F}_t$ 可測だが $\Prob_j$ に依存 & 残余請求権の価値 $V_j$ \\
\bottomrule
\end{tabularx}
\captionof{table}{観測できない量の三段階の主観性。}
\label{tab:subjectivity-levels}
\end{center}

$v$ の主観性は情報の非対称であり、開示によって解消しうる。
$V_j$ の主観性は\textbf{確率測度の不一致}であり、
同じ情報を見ても予測が異なるため開示では解消しない。

この区別から二つの帰結が出る。
第一に、$\kappa_K^{i}$ の値は当事者間で一致せず、他の $\kappa_i$ と性質が異なる。
第二に、\textbf{その不一致こそが出資を成立させる}。
出資者が経営者より $V$ を高く評価するから出資する。
これは命題\ref{prop:discount}の割引因子の異質性と同型の構造である。
\end{remark}

資本の期限については注\ref{rem:tau-K}で述べる。

\begin{remark}[$\tau_K$ と倒産]
\label{rem:tau-K}
$\tau_K$ は借入では返済期限として確定するが、出資では確定しない。
資本は最も長期の与信であり、期限を持たない。

配当を義務化すれば $\kappa_K$ が確定するが、
支払不能時に式\eqref{eq:constraint}が直ちに破れる。
義務化せず残余請求権とし、
第\ref{sec:objective}節の倒産時刻 $\tau$ において強制的に清算する仕組みが
これを回避している。破産は $\kappa_K$ の確定機構である。
\end{remark}

\begin{remark}[預り金の扱い]
$\kappa_i<0$ のうち、エスクローのように法的に他人に帰属する資金は、
式\eqref{eq:cash}上は同じ位置に現れるが性質が異なる。
流量が減少すると急速に流出するため、他の負の $\kappa$ と同一視してはならない。
\end{remark}

\chapter{情報構造と目的関数：契約が書ける条件と経路制約}
\label{sec:info}

\section{検証可能な共有情報}
\label{sec:verifiable}

主体 $i$ が観測できる情報を $\mathcal{G}_i\subseteq\mathcal{F}$ とし、
当事者間で検証可能な共有情報を $\mathcal{G}=\bigcap_i \mathcal{G}_i$ と書く。

\begin{proposition}[実装可能性]
\label{prop:implementable}
決済スケジュール $\pi$ が契約として実装可能であるための必要条件は、
$\pi$ が $\mathcal{G}$ 可測であることである。
\end{proposition}

\begin{proof}
$\pi$ が $\mathcal{G}$ 可測でないとする。
このとき $\mathcal{G}$ 上で区別できない二つの状態 $\omega,\omega'$ で
$\pi(\omega)\neq\pi(\omega')$ となるものが存在する。
支払義務が $\omega$ か $\omega'$ かに依存するにもかかわらず、
当事者はいずれであるかを共通に検証できない。
したがって支払額について合意に到達できず、第三者による履行強制もできない。
\end{proof}

この条件が $\Phi$ の設計に強い制約を与える。
努力水準 $a$ が $\mathcal{G}$ に含まれないとき、$a$ を条件とする契約は書けず、
観測可能な結果 $y$ を条件とするしかない。

\begin{center}
\begin{tabularx}{\textwidth}{@{}llX@{}}
\toprule
観測できるもの & 契約形態 & 例 \\
\midrule
努力 $a$ & 投入連動 & 時間給、コストプラス \\
結果 $y$ のみ & 成果連動 & 成果報酬、ロイヤルティ、免責 \\
どちらも部分的 & 混合 & 固定＋歩合 \\
\bottomrule
\end{tabularx}
\captionof{table}{観測可能なものと対応する契約形態。}
\label{tab:observability-contracts}
\end{center}

これは契約理論の標準的な結果であり（\cite{holmstrom1979}）、本稿の寄与ではない。

\begin{remark}[$\mathcal{G}$ が限るのは自由度(4) である]
\label{rem:G-limits-dof4}
命題\ref{prop:implementable}は、第\ref{sec:dof}節の自由度(4)（状態依存性）に
直接効く。$\pi$ を $\omega$ に依存させられるのは $\mathcal{G}$ の範囲内に限られ、
$\mathcal{G}$ の外にある $\omega$ の区別は契約に書けない。

\textbf{したがって $\mathcal{G}$ は $\Phi$ の定義域ではなく値域の側を限る。}
履行 $\delta$ は何でも選べるが、それに応じて $\pi$ をどう動かせるかが制約される。

この制約は第\ref{part:catalog}部の判定順序に現れる。
自由度(4) で判定される族は族5（成果・状態連動）であり、
$\pi=h(y)$ の $y$ が $\mathcal{G}$ 可測でなければこの族は実装できない。
第\ref{ch:protocol}章の族7の不可能性（命題\ref{prop:beneficiary}）も同じ形をとる。
\end{remark}

\begin{remark}[$\mathcal{G}$ の内生性]
\label{rem:G-endogenous}
本章は $\mathcal{G}$ を所与として扱う。実際には\textbf{事業者は $\mathcal{G}$ に投資できる}。
追跡調査を行えば結果 $y$ が観測可能になり、検証可能な情報の範囲が広がる。

第\ref{sec:jinzai-structure}節で扱う人材サービスの掲載項目に
「離職が判明せず」の人数欄が存在することは、
$\mathcal{G}$ の広さが事業者ごとに異なり、かつ選択の対象であることを示している。

したがって命題\ref{prop:implementable}は、
$\mathcal{G}$ を固定した短期の条件として読むべきである。
長期には $\mathcal{G}$ 自体が $\Phi$ の一部となる。
\end{remark}

\begin{remark}[方向の区別]
リスク移転には方向の異なる二種がある。
成果報酬とロイヤルティは、\emph{努力が観測できない}ためエージェント側にリスクを残す形式である。
一方コストプラスは、\emph{努力は観測できるが結果が不確実}なためプリンシパル側にリスクを戻す形式である。
両者を「リスク分担」として一括すると、この差が消える。
\end{remark}

\begin{example}[成果連動の最適強度]
\label{ex:bstar}
線形契約 $\pi = \alpha + b\,y$ を考える。
エージェントの絶対的リスク回避度を $\gamma$、結果のノイズ分散を $\sigma^2$ とすると、
（$r$ は第\ref{sec:growth}章で売上速度に用いるため、ここでは $\gamma$ と書く）
最適な成果連動係数は
\begin{equation}
b^{\star} = \frac{1}{1+\gamma\,\sigma^{2}\,k}
\label{eq:bstar}
\end{equation}
の形をとる（$k$ は努力の限界費用に関する定数、\cite{holmstrom1991}）。
$\sigma^2$ が大きいほど $b^\star$ は小さい。
すなわち\textbf{結果が騒がしいほど成果連動は弱めるべきである}。

数値例として $\gamma k = 1$ とすると、$\sigma^2=0.5$ で $b^\star=0.67$、
$\sigma^2=2$ で $b^\star=0.33$、$\sigma^2=9$ で $b^\star=0.10$ となる。
\end{example}

\section{目的関数：経路ごとに課される現金制約}
\label{sec:objective}

制約に現れる $M(t,\omega)$ は命題\ref{prop:cash}の現金残高である。

\begin{definition}[最適化問題]
期間利益を $\phi_t$、割引因子を $\beta$ として、
\begin{align}
\max_{\Phi}\quad & \E\!\left[\sum_{t\in T}\beta^{t}\phi_t\right] \label{eq:obj}\\
\text{s.t.}\quad & M(t,\omega)\geq 0 \qquad \forall t\in T,\ \forall\omega\in\Omega \label{eq:constraint}
\end{align}
\end{definition}

式\eqref{eq:constraint}の性質が本稿の中心である。

\begin{proposition}[経路制約の非可換性]
\label{prop:pathwise}
$\E[M(t)]\geq 0$ は式\eqref{eq:constraint}を含意しない。
また倒産時刻 $\Tins=\inf\{t: M(t)<0\}$ は吸収状態であり、
$t>\Tins$ における $\phi_t$ は実現しない。
$\tau_i$ は決済の遅れに用いているため、倒産時刻は $\Tins$ と書く。
\end{proposition}

\begin{proof}
第一の主張は反例で足りる。
$M(t)$ が確率 $1/2$ で $+100$、確率 $1/2$ で $-100$ をとるとき
$\E[M(t)]=0\geq 0$ であるが、確率 $1/2$ の経路で制約が破れる。

第二の主張は、$M(\Tins)<0$ において支払不能が生じ、
以後の履行 $\delta$ が実行できないことによる。
したがって目的関数は正しくは
$\E\bigl[\sum_{t<\Tins}\beta^t\phi_t\bigr]$ と書かれるべきであり、
$\Tins$ は $\Phi$ に依存する。
\end{proof}

\begin{corollary}[黒字倒産]
\label{cor:insolvency-with-profit}
$\phi_t>0$ が全ての $t$ で成立していても、
式\eqref{eq:constraint}が破れうる。すなわち黒字倒産は
目的関数\eqref{eq:obj}の失敗ではなく制約\eqref{eq:constraint}の失敗である。
\end{corollary}

\begin{proof}
第\ref{sec:growth}章の式\eqref{eq:cash-net}より $M=E-W-A$ である。
$\phi>0$ は式\eqref{eq:equity}を通じて $E$ を増やすが、
$W$ の増加がそれを上回れば $M$ は減少する。
命題\ref{prop:depletion}はこの場合の枯渇時刻を与える。
\end{proof}

黒字倒産の条件は第\ref{sec:growth}章で $g>g^\star$ として書き直される（注\ref{rem:level-and-change}）。

\begin{example}[黒字倒産の数値例]
\label{ex:insolvency}
粗利率 $m=10\%$ の事業が、年商1億円から2億円へ成長する。利益は 2,000万円である。

顧客への与信が売上の 20\%、在庫が売上の 5\% で推移するとすれば、
式\eqref{eq:kappa}の信用ポジションと在庫の合計は
1億円のとき 2,500万円、2億円のとき 5,000万円となる。
増加分 2,500万円が現金から出ていく。

利益 2,000万円では 500万円不足する。
損益は黒字、現金は不足という状態であり、式\eqref{eq:constraint}が破れうる。
\end{example}

\section{割引因子の異質性}

\begin{proposition}[決済スケジュールの組み替え利得]
\label{prop:discount}
$\beta_C<\beta_0$（顧客が企業より将来を強く割り引く）であるとき、
履行 $\delta$ を変えずに $\pi$ の時間配置のみを変更することで、
両者の主観的評価を同時に改善する組み替えが存在する。
\end{proposition}

\begin{proof}
現行契約の支払いを時点 $0$ における $p$ とする。
これを時点 $1$ の $p'$ に繰り延べる変更を考える。
顧客の評価は $-p$ から $-\beta_C p'$ に変わるから、
$\beta_C p' < p$ すなわち $p' < p/\beta_C$ であれば顧客は改善する。
企業の評価は $p$ から $\beta_0 p'$ に変わるから、
$\beta_0 p' > p$ すなわち $p' > p/\beta_0$ であれば企業も改善する。
$\beta_C<\beta_0$ より $p/\beta_0 < p/\beta_C$ であり、
この区間は空でない。任意の $p'$ をこの区間に取れば両者が改善する。
\end{proof}

\begin{example}[割賦の利得]
年利換算で顧客の割引率が 20\%（$\beta_C=0.833$）、
企業の資金コストが 5\%（$\beta_0=0.952$）であるとする。
現金価格 100,000円の商品について、
1年後払いの価格 $p'$ は $105{,}000 < p' < 120{,}000$ の範囲で
双方が改善する。$p'=112{,}000$ とすれば、
企業は割引現在価値で $112{,}000\times 0.952 = 106{,}624$ 円を得て 6,624円の増益、
顧客は主観的に $112{,}000\times 0.833=93{,}296$ 円と評価し 6,704円の得となる。
\end{example}

\begin{remark}[解釈の分岐]
$\beta_C<\beta_0$ には二つの異なる原因がありうる。
流動性制約に直面している顧客にとって、現在の資金は客観的に価値が高い。
一方、双曲割引（\cite{laibson1997}）などの認知的性質による場合もある。
前者に応えることと後者を利用することは、外形上ほぼ区別できないが、
耐久性（第\ref{sec:surplus}章）は正反対である。この区別は本枠組みの内部では決着しない。
\end{remark}

\chapter{成長制約：\texorpdfstring{$\CCC$}{CCC} と自己金融可能成長率}
\label{sec:growth}

\section{キャッシュコンバージョンサイクル}
\label{sec:ccc-def}

売上速度を $r(t)$、運転資本を定義\ref{def:wc}の $W(t)$ とする。
$W$ は式\eqref{eq:kappa}の $\kappa_i$ の和であるから、$\Phi$ に依存する。

\begin{definition}[キャッシュコンバージョンサイクル]
定常状態において
\begin{equation}
\CCC = \frac{W}{r}
\label{eq:ccc}
\end{equation}
と定義する。次元は時間である。
\end{definition}

\section{基本不等式}
\label{sec:fundamental}

命題\ref{prop:cash-net}の式\eqref{eq:cash-net}に式\eqref{eq:ccc}を代入すると、現金残高が
\begin{equation}
M(t) = E(t) - \CCC\cdot r(t) - A(t)
\label{eq:M-ccc}
\end{equation}
と書ける。第\ref{sec:objective}節の制約 $M\geq 0$ はしたがって次の形をとる。

\begin{proposition}[基本不等式]
\label{prop:fundamental}
式\eqref{eq:constraint}の現金制約は
\begin{equation}
\CCC\cdot r(t,\omega) \ \leq\ E(t) - A(t)
\label{eq:fundamental}
\end{equation}
と同値である。
\end{proposition}

\begin{proof}
式\eqref{eq:cash-net}に $W=\CCC\cdot r$ を代入して $M\geq 0$ を整理すればよい。
\end{proof}

式\eqref{eq:fundamental}は\textbf{運転資本が自己資本から固定資産を引いた額を超えられない}
と読む。左辺は $\Phi$ が決め、右辺は資本構成が決める。
本章の残りと、第\ref{part:catalog}部の各章は、この不等式の特例を扱う。

\begin{corollary}[規模の上限]
\label{cor:rmax}
$\CCC>0$ のとき、実行可能な売上速度は
\begin{equation}
r \ \leq\ r^{\max} = \frac{E-A}{\CCC}
\label{eq:rmax}
\end{equation}
で上から抑えられる。$\CCC\leq 0$ かつ $E\geq A$ のとき、
式\eqref{eq:fundamental}は任意の $r$ で成立する。
\end{corollary}

次元は円÷年であり、売上速度と一致する。
$g^\star$ が成長率の上限であるのに対し、$r^{\max}$ は\textbf{水準の上限}である。
両者は別の制約であり、一方は他方を含意しない（注\ref{rem:level-and-change}）。

\begin{corollary}[資本がなければ $\CCC\leq 0$ の $\Phi$ しか選べない]
\label{cor:zero-equity}
$E=0$ のとき、式\eqref{eq:fundamental}は $\CCC\cdot r\leq -A\leq 0$ を要求する。
すなわち $\CCC\leq 0$ である。
\end{corollary}

系\ref{cor:zero-equity}は第\ref{ch:solo}章の出発点になる。
資本を持たない主体に残るのは、決済が履行に先行する $\Phi$ だけである。

式\eqref{eq:ccc}の $\CCC$ を用いる。

\section{成長とフリーキャッシュフロー}
\label{sec:fcf}

基本不等式は水準についての制約である。次に変化を扱う。

\begin{proposition}[成長とフリーキャッシュフロー]
\label{prop:fcf}
粗利率を $m$、設備投資を $I$ とし、
\textbf{$\CCC$ を一定かつ $m$ と独立とすると}
\begin{equation}
\FCF = m\,r - \CCC\cdot\dot{r} - I
\label{eq:fcf}
\end{equation}
が成立する。
\end{proposition}

\begin{proof}
フリーキャッシュフローは、粗利から運転資本の増加分と設備投資を差し引いたものである。
\[
\FCF = m\,r - \dot{W} - I .
\]
式\eqref{eq:ccc}より $W=\CCC\cdot r$ であり、$\CCC$ が一定であれば
$\dot W = \CCC\cdot\dot r$ を得る。代入すれば式\eqref{eq:fcf}が従う。
\end{proof}

\begin{remark}[$\FCF$ は現金残高の時間微分である]
\label{rem:fcf-is-mdot}
式\eqref{eq:M-ccc}を微分する。式\eqref{eq:equity}より $\dot E=\phi-\dot{\mathrm{Div}}$、
固定資産は投資から減価償却を引いて $\dot A = I - d\,r$ である。
配当と外部調達がなければ
\[
\dot M = \phi + d\,r - \CCC\,\dot r - I = \FCF
\]
となる。$\FCF$ は独立に導入された量ではなく、\textbf{式\eqref{eq:cash-net}の時間微分である}。
したがって次節の $g^\star$ は $\dot M\geq 0$ の境界にあたる。
\end{remark}

式\eqref{eq:fcf}を $g$ について解く。分母には式\eqref{eq:ccc}の $\CCC$ が現れる。

\begin{corollary}[自己金融可能成長率]
\label{cor:gstar}
売上高減価償却費率を $d$ とし、成長率を $g=\dot r/r$ とすると、
$\FCF\geq 0$ の条件は
% 本稿で唯一、枠で囲む式。PDF は \boxed でよいが、HTML では tex4ht が
% <menclose notation="box"> を出す。これは MathML Core に無いため
% Chrome と Safari では枠が描かれず、Firefox では枠の上辺が分数に切られる。
% HTML 側は空の span で印をつけ、隣接する式の表に CSS で枠を描く。
% 生の <div> を段落の途中に差し込むと <p> の入れ子が壊れ、make4ht の
% domfilter が XML 解析を諦めて HTML 解析に退避する。そのときファイル全体が
% 再直列化され、数式中の < が <mo>&lt;</mo> から壊れた出力に化ける。
% 印を釣り合う空要素にすればこれが起きない。
\ifdefined\HCode
\HCode{<span class="keyeq-mark"></span>}%
\begin{equation}
g \ \leq\ g^{\star} = \frac{m + d - I/r}{\CCC}
\label{eq:gstar}
\end{equation}
\else
\begin{equation}
\boxed{\ g \ \leq\ g^{\star} = \frac{m + d - I/r}{\CCC}\ }
\label{eq:gstar}
\end{equation}
\fi
である。$\CCC<0$ のとき $g^\star$ は存在せず、任意の成長率で $\FCF>0$ となる。
\end{corollary}

\begin{proof}
$m$ を営業利益率とすると減価償却費が控除済みであるが、
減価償却は現金支出を伴わない。したがって現金ベースの粗利は $(m+d)\,r$ である。
式\eqref{eq:fcf}を修正すると
$\FCF = (m+d)\,r - \CCC\,\dot r - I$。
$\dot r = g\,r$ を代入して $r$ で割れば
$\FCF/r = (m+d) - \CCC\,g - I/r$。
$\FCF\geq 0$ を $g$ について解けば式\eqref{eq:gstar}を得る。
$\CCC<0$ のときは左辺が全ての $g>0$ で正となる。
\end{proof}

\section{水準と変化は別の制約である}
\label{sec:level-and-change}

\begin{remark}[水準の制約と変化の制約は独立である]
\label{rem:level-and-change}
系\ref{cor:rmax}は $M\geq 0$ から、系\ref{cor:gstar}は $\dot M\geq 0$ から出る。
一方は他方を含意しない。$M(0)>0$ であっても $g>g^\star$ ならば $M$ は減少し、
有限の時刻で式\eqref{eq:constraint}が破れる。
\end{remark}

\begin{proposition}[現金の枯渇時刻]
\label{prop:depletion}
$I=0$、$g$ 一定、$r(t)=r_0e^{gt}$ とする。$g>g^\star=(m+d)/\CCC$ のとき、
$M(0)=M_0>0$ から出発しても現金は
\begin{equation}
\Tins = \frac{1}{g}\,
\ln\!\left(1+\frac{g\,M_0}{(\CCC\,g-m-d)\,r_0}\right)
\label{eq:depletion}
\end{equation}
で尽きる。$g\downarrow g^\star$ のとき $\Tins\to\infty$ である。
\end{proposition}

\begin{proof}
注\ref{rem:fcf-is-mdot}より $\dot M = a\,r_0e^{gt}$、$a=(m+d)-\CCC g$ である。
積分して $M(t)=M_0+(a r_0/g)(e^{gt}-1)$。
$g>g^\star$ より $a<0$ であるから $M$ は減少し、$M(\Tins)=0$ を解けば
式\eqref{eq:depletion}を得る。$g\to g^\star$ で $a\to 0$ となり対数の引数が発散する。
\end{proof}

式\eqref{eq:depletion}は、命題\ref{prop:pathwise}が記号で導入した倒産時刻に値を与え、
系\ref{cor:insolvency-with-profit}の黒字倒産がいつ現実化するかを定める。

\begin{example}[$g^\star$ は閾値だが崖ではない]
\label{ex:depletion}
$m=5\%$、$d=2\%$、$\CCC=0.25$ 年、$M_0=100$、$r_0=1{,}000$ とすると $g^\star=28\%$ である。
\begin{center}
\begin{tabular}{@{}lr@{}}
\toprule
成長率 & $\Tins$ \\
\midrule
$g^\star+0.001$ & 16.8 年 \\
$g^\star+0.01$ & 8.7 年 \\
$g^\star+0.05$ & 3.9 年 \\
$g^\star+0.2$ & 1.4 年 \\
$g^\star+0.5$ & 0.6 年 \\
\bottomrule
\end{tabular}
\captionof{table}{$g^\star$ の超過幅と現金の枯渇時刻。}
\label{tab:depletion}
\end{center}
わずかな超過には十数年の猶予があり、大幅な超過は一年で尽きる。
\textbf{$g^\star$ を超えたかどうかではなく、どれだけ超えたかが効く。}
\end{example}

\begin{remark}[$\CCC$ は需要変動の増幅係数でもある]
\label{rem:ccc-amplifier}
式\eqref{eq:M-ccc}は $r$ について線形であり、係数は $-\CCC$ である。
$E$ と $A$ が短期に $r$ と独立であれば
\[
\operatorname{Var}[M] = \CCC^{2}\operatorname{Var}[r]
\]
となる。\textbf{$\CCC$ は水準の上限を決めるだけでなく、
需要の変動を現金の変動に増幅する。}
$\CCC$ が大きい業種ほど、同じ需要変動に対して式\eqref{eq:constraint}が破れやすい。

本稿は $r$ の過程を指定しないため $\Prob(\Tins\leq t)$ を与えない。
増幅係数が $\CCC$ であることは過程の指定なしに言える。
\end{remark}

\begin{remark}[$I \approx d\,r$ の仮定を置かない]
\label{rem:no-steady-investment}
資産規模が一定であれば $I = d\,r$ となり、式\eqref{eq:gstar}は $m/\CCC$ に帰着する。
本稿はこの仮定を置かない。第\ref{sec:gstar-measured}節の測定によれば、
$I/r$ は業種と規模によって $-1.8\%$ から $6.7\%$ まで分布し、
$d$ と一致しない。
\end{remark}

\begin{remark}[$\CCC$ が小さい場合の発散]
\label{rem:gstar-divergence}
式\eqref{eq:gstar}は $\CCC\to 0$ で発散する。
$\CCC$ が正であっても数日程度であれば、分子の微小な変動が
$g^\star$ を大きく動かす。

たとえば $\CCC=5$ 日のとき、分子が 1 ポイント変化すると
$g^\star$ は約 73 ポイント動く。
\textbf{$\CCC$ が小さい業種では $g^\star$ は指標として機能しない。}
第\ref{sec:gstar-measured}節ではこれを標準偏差の併記により示す。
\end{remark}

$g^\star$ は自己金融可能な成長率であり、\cite{churchill2001}の
self-financeable growth rate と実質的に同一である。
$\FCF$ と成長率の関係を図\ref{fig:growth}に示す。

\begin{figure}[htbp]
\centering
\insertfig{growth-fcf}
\caption{成長率とフリーキャッシュフロー。縦軸切片が式\eqref{eq:gstar}の分子、$\CCC$ の符号が傾きの符号を決める。}
\label{fig:growth}
\end{figure}

\begin{example}[$g^\star$ の実際の値]
営業利益率 $m=5\%$、減価償却費率 $d=2.5\%$、設備投資 $I/r=3.5\%$ の事業を考える。
第\ref{sec:gstar-measured}節で測定した全産業の値に近い。
分子は
\[
m + d - I/r = 5.0 + 2.5 - 3.5 = 4.0\%
\]
となる。$\CCC$ の値ごとに $g^\star$ を計算する。

\begin{center}
\begin{tabular}{@{}rrrr@{}}
\toprule
$\CCC$（日） & $\CCC$（年） & $g^\star$ & 参考：$m/\CCC$ \\
\midrule
120 & 0.329 & 12.2\% & 15.2\% \\
90 & 0.247 & 16.2\% & 20.3\% \\
60 & 0.164 & 24.3\% & 30.4\% \\
30 & 0.082 & 48.7\% & 60.9\% \\
\textminus 30 & \textminus 0.082 & 存在しない & 存在しない \\
\bottomrule
\end{tabular}
\captionof{table}{$\CCC$ の値ごとの $g^\star$ と $m/\CCC$（数値例）。}
\label{tab:gstar-example}
\end{center}

$\CCC=120$ 日の事業は年率12\%を超える成長で必ず外部資金を要する。
一方 $\CCC=-30$ 日の事業は、成長が速いほど現金が増える。
\textbf{符号が定性的な体質を、絶対値が定量的な余裕を決める。}

右端の列は $d$ と $I$ を無視した場合の値である。
本例では $d<I/r$ であるため $m/\CCC$ は $g^\star$ を過大に評価する。
\textbf{減価償却を上回る投資を行う事業では、$m$ のみによる評価は楽観的になる。}
\end{example}

\begin{remark}[定義域]
式\eqref{eq:ccc}は定常状態を前提としており、$r\approx 0$ の創業直後には定義されない。
この点は第\ref{sec:age}節で改めて扱う。
\end{remark}

\section{\texorpdfstring{$\CCC$}{CCC} は遅れ時間の和である}
\label{sec:ccc-as-tau}

命題\ref{prop:kappa-tau}を式\eqref{eq:ccc}に代入すると、$r$ が消える。

\begin{proposition}[$\CCC$ の構造]
\label{prop:ccc-structure}
定常状態において
\begin{equation}
\CCC = \sum_i \tau_i + \tau_{\mathrm{inv}},
\qquad \tau_{\mathrm{inv}} \equiv \frac{\mathrm{Inv}}{r}
\label{eq:ccc-tau}
\end{equation}
が成立する。すなわち $\CCC$ は流量に依存せず、遅れ時間の和である。
\end{proposition}

\begin{proof}
$W=\sum_i\kappa_i+\mathrm{Inv}$ と式\eqref{eq:kappa-tau}より
$W = r\sum_i\tau_i + \mathrm{Inv}$。両辺を $r$ で割れば従う。
なお式\eqref{eq:kappa-tau}は時間平均として成立する（注\ref{rem:kappa-tau-average}）ため、
本命題も平均についての主張である。
\end{proof}

式\eqref{eq:ccc-tau}から三つの帰結が出る。いずれも演繹であり、観測を用いない。

\begin{corollary}[水準差の源泉]
\label{cor:ccc-level}
$\tau_i$ は第\ref{sec:dof}節の自由度(1)、すなわち
$\operatorname{supp}(\pi)$ と $\operatorname{supp}(\delta)$ の前後関係が指定する量である。
したがって $\CCC$ の水準は $\Phi$ が決める。
\end{corollary}

\begin{corollary}[横断比較の無意味さ]
\label{cor:no-cross-comparison}
異なる $\Phi$ を持つ対象の間で $\CCC$ の水準を比較することは、
異なる契約形態を比較することであり、
需要や効率についての情報を持たない。
\end{corollary}

\begin{corollary}[集計は $r$ 加重平均になる]
\label{cor:ccc-weighted}
命題\ref{prop:kappa-additive}の並置において、部分ごとの売上速度を $r_k$ とすると
\begin{equation}
\CCC = \sum_k w_k\,\CCC_k, \qquad w_k = \frac{r_k}{\sum_j r_j}
\label{eq:ccc-weighted}
\end{equation}
であり、$\CCC$ は加法的でない。
\end{corollary}

\begin{proof}
$W=\sum_k W_k$（命題\ref{prop:kappa-additive}）と $W_k=\CCC_k r_k$ より
$\CCC=\sum_k \CCC_k r_k/\sum_j r_j$ を得る。
\end{proof}

式\eqref{eq:ccc-weighted}は、異なる族の $\Phi$ を並行して運用する企業について
$\CCC$ を計算すると異符号の加重平均を観測することを意味する。
観測単位を企業ではなく $\Phi$ に置く根拠がここにある。

\begin{corollary}[変化の分解]
\label{cor:ccc-change}
式\eqref{eq:ccc-tau}の差分は
\begin{equation}
\Delta\CCC = \underbrace{\sum_i \Delta\tau_i}_{\text{契約の変化}}
 + \underbrace{\Delta\tau_{\mathrm{inv}}}_{\text{運用の変化}}
\label{eq:ccc-change}
\end{equation}
と分かれる。式\eqref{eq:ccc-change}において $\Delta\tau_i\neq 0$ は契約条件の変更を意味し、
$\Delta\tau_{\mathrm{inv}}$ は契約と独立に生じる。
\end{corollary}

\section{経路制約と尖頭}
\label{sec:peak}

式\eqref{eq:constraint}はすべての $t$ と $\omega$ について課される。
一方、注\ref{rem:kappa-tau-average}のとおり $\bar\kappa_i=r\tau_i$ は時間平均である。
\textbf{拘束するのは平均ではなく尖頭である。}

\begin{proposition}[尖頭の運転資本]
\label{prop:peak}
売上速度が一定、締め期間 $T_b$、支払サイト $\tau_i$ のとき、
$\kappa_i(t)$ の時間平均は $r(\tau_i+T_b/2)$、最大値は $r(\tau_i+T_b)$ である。
したがって式\eqref{eq:fundamental}は
\begin{equation}
r \ \leq\ \frac{E-A}{\CCC + T_b/2}
\label{eq:rmax-peak}
\end{equation}
となる。
\end{proposition}

\begin{proof}
締め時点 $kT_b$ に $[(k-1)T_b,kT_b)$ の履行が請求され、$kT_b+\tau_i$ に決済される。
未決済の累積は請求前の経過分と請求済み未決済分の和であり、
締め直後に最小 $r\tau_i$、締め直前に最大 $r(\tau_i+T_b)$ をとる鋸歯となる。
一周期にわたる平均は $r(\tau_i+T_b/2)$ である。
平均で測った $\CCC$ に対し尖頭は $T_b/2$ だけ大きいから、
式\eqref{eq:fundamental}の左辺を尖頭で評価すれば式\eqref{eq:rmax-peak}を得る。
\end{proof}

月次締めであれば $T_b/2$ は約15日である。
\textbf{平均から計算した $\CCC$ で必要な資本を見積もると、半月分だけ過小になる。}
第\ref{part:results}部の測定は期末残高または期中平均に基づいており、この差を含まない。

\section{比較静学}
\label{sec:comparative-statics}

本章の量はいずれもパラメータの関数である。
微分を並べると、$\Phi$ の設計者が何をどれだけ動かせばよいかが読める。

\begin{center}
\begin{tabularx}{\textwidth}{@{}llX@{}}
\toprule
量 & 微分 & 内容 \\
\midrule
$g^\star$ & $\partial\ln g^\star/\partial\tau_i = -1/\CCC$ & $i$ に依存しない \\
$g^\star$ & $\partial g^\star/\partial m = 1/\CCC$ & 利益率の改善は $\CCC$ が小さいほど効く \\
式\eqref{eq:rmax} & $\partial\ln r^{\max}/\partial\tau_i = -1/\CCC$ & $g^\star$ と同一の係数 \\
式\eqref{eq:rmax} & $\partial\ln r^{\max}/\partial E = 1/(E-A)$ & 増資の効果は $A$ が大きいほど強い \\
$\operatorname{Var}[M]$ & $\partial\operatorname{Var}[M]/\partial\CCC = 2\CCC\operatorname{Var}[r]$ & 注\ref{rem:ccc-amplifier} \\
\bottomrule
\end{tabularx}
\captionof{table}{本章の量の比較静学。}
\label{tab:comparative-statics}
\end{center}

\begin{proposition}[遅れ時間は限界で完全代替である]
\label{prop:tau-substitution}
式\eqref{eq:ccc-tau}より、任意の $i$ について
$\partial\ln g^\star/\partial\tau_i = -1/\CCC$ であり、$i$ に依存しない。
\end{proposition}

\begin{proof}
$g^\star=(m+d-I/r)/\CCC$ と $\CCC=\sum_j\tau_j+\tau_{\mathrm{inv}}$ より
$\partial\CCC/\partial\tau_i=1$、したがって
$\partial\ln g^\star/\partial\tau_i = -\partial\ln\CCC/\partial\tau_i = -1/\CCC$ を得る。
\end{proof}

\textbf{回収を一日早めることと、支払を一日延ばすことは、$g^\star$ に対して等価である。}
表\ref{tab:comparative-statics}の第三行より、$r^{\max}$ についても同じ係数が現れる。
水準と変化のどちらを見ても、$\Phi$ の設計指針は変わらない。

系\ref{cor:no-cross-comparison}は、
注\ref{rem:level-vs-change}で述べた制約に演繹的な根拠を与える。
本稿の量が水準として定義されていることは記述上の選択だが、
\textbf{横断比較ができないことは定義からの帰結}である。

\begin{remark}[比較可能な量を構成できない]
\label{rem:no-normalization}
系\ref{cor:no-cross-comparison}は比較できない理由を述べるが、
比較する方法を与えない。

自然な候補は、契約が指定する $\tau_i^{\mathrm{contract}}$ からの乖離
\[
\tau_i^{\mathrm{observed}} - \tau_i^{\mathrm{contract}}
\]
である。この量は $\Phi$ に依存しないため業種を跨いで比較でき、
契約どおりに履行されているかを測る。

しかし $\tau_i^{\mathrm{contract}}$ は個別契約の支払条件であり、公表統計に存在しない。
\textbf{理論の問題は解けたが、測定の問題に移った}
（第\ref{sec:obstacle-taxonomy}節の分類では（iv）アクセス制約）。
\end{remark}

\begin{remark}[$m$ と $\CCC$ の独立性]
\label{rem:m-ccc-independence}
系\ref{cor:gstar}は $g^\star$ を比として与えるが、
分子と分母が独立に動くことを保証しない。

第\ref{sec:gstar-measured}節の測定では、2000年度から2024年度にかけて
$\CCC$ が 49\%、$m$ が 91\% と\textbf{ともに上昇している}。
運転資本を厚く持つことと利益率が高いことが正の関係にあるならば、
$g^\star$ の変化を「制約の緩和」と解釈することはできない。

本稿は両者の関係を定式化しない。
$g^\star$ は比の値として報告するにとどめ、因果的な解釈を与えない。
\end{remark}

\chapter{余剰とレイヤー：源泉・耐久性・可変性}
\label{sec:surplus}

\section{余剰の三分解}
\label{sec:surplus-decomp}

注\ref{rem:Phi-vs-phi}で予告した余剰 $\phi$ を、ここで定める。
期間利益を源泉によって三分解する。この分解は本稿の整理である。
第三項は式\eqref{eq:cog}の認知余剰であり、
式\eqref{eq:phi}の自由度(5) が立つときに正となりうる。
第一項と第二項は定義\ref{def:phi-prod}で定める。
\begin{equation}
\phi = \phi_{\mathrm{prod}} + \phi_{\mathrm{barg}} + \phi_{\mathrm{cog}}
\label{eq:surplus}
\end{equation}

\begin{center}
\begin{tabularx}{\textwidth}{@{}llX@{}}
\toprule
記号 & 源泉 & 性質と減衰 \\
\midrule
$\phi_{\mathrm{prod}}$ & 限界費用 $c(\delta)$ の優位 & 技術的優位。他者の取り分を減らさない。競争で削られるが投資で再生産できる \\
$\phi_{\mathrm{barg}}$ & 写像 $\Phi$ & 取り分の移転。$\sum_i \phi_i = 0$ のゼロサム。制度層が制約を課す \\
$\phi_{\mathrm{cog}}$ & $\dbar$ と $\delta$ の乖離 & 学習・規制・UI改善で削られ、原則として回復しない \\
\bottomrule
\end{tabularx}
\captionof{table}{余剰 $\phi$ の三分解と、各成分の源泉・耐久性。}
\label{tab:surplus-decomposition}
\end{center}

\begin{remark}[リスク引受の対価はどこに入るか]
\label{rem:phi-risk}
第\ref{sec:three-ops}節の移転と集約を業として行う主体、
すなわち保険者や決済代行のような仲介者は、
リスクを引き受けた対価を得る。これは式\eqref{eq:surplus}のどこに属するか。

集約によって実際に分散が減少した部分は $\phi_{\mathrm{prod}}$ である。
命題\ref{prop:pooling}により $n$ を増やせば $\sigma^2/n$ の項は消え、
これは技術的な優位であって他者の取り分を減らさない。

一方、式\eqref{eq:pooling}の極限 $\rho\sigma^2$ に対応する部分は消せない。
これを引き受ける対価は移転であり、$\sum_i\phi_i=0$ を満たす。
すなわち $\phi_{\mathrm{barg}}$ に属する。

実務上、両者は分離されずに一つの料率として現れる。
本稿では便宜的に $\phi_{\mathrm{risk}}$ と表記する場合があるが
（第\ref{sec:pred-J}節）、これは新たな成分ではなく、
$\phi_{\mathrm{prod}}$ と $\phi_{\mathrm{barg}}$ の未分離な和である。
\end{remark}

\begin{definition}[生産余剰と交渉余剰]
\label{def:phi-prod}
定義\ref{def:cost}の限界費用 $c(\delta)$ と、
支配力が働かない場合の価格 $p^{\ast}$ を用いて
\begin{align}
\phi_{\mathrm{prod}} &= \bigl(p^{\ast} - c(\delta)\bigr)\,\delta \\
\phi_{\mathrm{barg}} &= \bigl(p - p^{\ast}\bigr)\,\delta
\end{align}
と定める。$p$ は実際の価格である。$p^{\ast}$ の性質は注\ref{rem:counterfactual-price}で述べる。
\end{definition}

\begin{remark}[$p^{\ast}$ は反実仮想である]
\label{rem:counterfactual-price}
定義\ref{def:phi-prod}は $p^{\ast}$ を必要とするが、これは観測されない。
支配力が働かない状態の価格を想定する反実仮想であり、
産業組織論で市場支配力を測る際の標準的な手続きと同じ構造を持つ。

したがって $\phi_{\mathrm{prod}}$ と $\phi_{\mathrm{barg}}$ の分離には、
$p^{\ast}$ を特定する追加の仮定が要る。

要る仮定は一つである。\textbf{支配力以外の条件が等しい対照が観測できること。}
これがあれば対照の価格を $p^{\ast}$ と読み、差を $\phi_{\mathrm{barg}}$ に帰せる。
第\ref{sec:platform-fee}章はこの仮定を、
同一プラットフォーム内で機能のみが異なる二つの料率という形で満たす。
仮定が満たされない場合、本稿は $\phi_{\mathrm{prod}}+\phi_{\mathrm{barg}}$ の
和までしか述べない。
\end{remark}

認知余剰の減衰を、取引頻度 $\nu$ の関数として
\begin{equation}
\dot{\phi}_{\mathrm{cog}} = -\lambda(\nu)\,\phi_{\mathrm{cog}},
\qquad \lambda'(\nu)>0
\label{eq:decay}
\end{equation}
とモデル化する。

\begin{proposition}[減衰の含意]
式\eqref{eq:decay}の解は $\phi_{\mathrm{cog}}(t)=\phi_{\mathrm{cog}}(0)e^{-\lambda(\nu)t}$ であり、
半減期は $\ln 2/\lambda(\nu)$ である。
すなわち取引頻度が高いほど認知余剰は速く消える。
\end{proposition}

\begin{proof}
線形一階常微分方程式の直接積分による。
$\phi(t)=\phi(0)e^{-\lambda t}=\phi(0)/2$ を解けば $t=\ln 2/\lambda$ を得る。
$\lambda'(\nu)>0$ より半減期は $\nu$ の減少関数である。
\end{proof}

\begin{example}[頻度と半減期]
$\lambda$ が年間取引回数に比例し、$\lambda = 0.02\nu$ とする。
\begin{center}
\begin{tabular}{@{}llrr@{}}
\toprule
財 & 年間取引回数 $\nu$ & $\lambda$ & 半減期 \\
\midrule
日用品 & 200 & 4.0 & 0.17 年 \\
月額サービス & 12 & 0.24 & 2.9 年 \\
年契約（保険、会員権） & 1 & 0.02 & 34.7 年 \\
\bottomrule
\end{tabular}
\captionof{table}{取引頻度 $\nu$ と認知余剰の半減期（数値例）。}
\label{tab:decay-halflife-example}
\end{center}

\section{余剰と現金の関係}
\label{sec:surplus-cash}

$\phi$ は式\eqref{eq:equity}を通じてのみ現金に効く。三成分の効き方は同じではない。

\begin{proposition}[認知余剰の認識は現金を動かさない]
\label{prop:cog-cash-neutral}
前払い契約で $p\,\dbar$ を受領し、履行が $\delta<\dbar$ で終わったとする。
期限到来により残債務が消滅して $\phi_{\mathrm{cog}}$ が認識されるとき、
式\eqref{eq:cash-net}の $M$ は変化しない。
\end{proposition}

\begin{proof}
評価を $v(\delta)=p\,\delta$ とすると、認識前の信用ポジションは
$\kappa_C = p\delta - p\dbar = -p(\dbar-\delta) = -\phi_{\mathrm{cog}}$ である。
認識により $\kappa_C$ は $0$ になり、$W$ は $\phi_{\mathrm{cog}}$ だけ増える。
同時に式\eqref{eq:equity}より $E$ が同額だけ増える。
式\eqref{eq:cash-net}において両者は打ち消し合う。
\end{proof}

\textbf{認知余剰は、認識の時点では現金を生まない。}
現金は前払いを受け取った時点で既に入っており、
そこでの効果は $\phi_{\mathrm{cog}}$ ではなく $\kappa_C<0$ として現れていた。
式\eqref{eq:cog}の $\phi_{\mathrm{cog}}$ と、
第\ref{sec:kappa}章の未履行残高は、同じ取引の別の断面である。

\begin{center}
\begin{tabularx}{\textwidth}{@{}lX@{}}
\toprule
成分 & 現金が動く時点 \\
\midrule
$\phi_{\mathrm{prod}}$、$\phi_{\mathrm{barg}}$ & 決済時点。$\tau_C$ だけ遅れる \\
$\phi_{\mathrm{cog}}$ & 契約時点。認識時点ではない \\
\bottomrule
\end{tabularx}
\captionof{table}{式\eqref{eq:surplus}の三成分と現金が動く時点。}
\label{tab:surplus-cash-timing}
\end{center}

\begin{remark}[会計上の扱いと一致する]
命題\ref{prop:cog-cash-neutral}は、収益認識基準における
契約負債からの振替（breakage）と同じ内容である。
本稿における意味は、この振替が式\eqref{eq:cash-net}の $E$ と $W$ の
両方を同額動かすため相殺される、という形で $\kappa$ の言葉に翻訳される点にある。
\end{remark}

\begin{corollary}[認知余剰から積める自己資本には上限がある]
\label{cor:cog-bounded}
式\eqref{eq:decay}に従って減衰するとき、
$\phi_{\mathrm{cog}}$ から蓄積できる自己資本の総額は
\begin{equation}
\int_0^{\infty}\phi_{\mathrm{cog}}(0)\,e^{-\lambda(\nu)s}\,ds
= \frac{\phi_{\mathrm{cog}}(0)}{\lambda(\nu)}
\label{eq:cog-total}
\end{equation}
で抑えられる。$\lambda'(\nu)>0$ より、取引頻度が高いほど上限は小さい。
\end{corollary}

式\eqref{eq:cog-total}を系\ref{cor:rmax}に代入すると、認知余剰以外の源泉がない事業について
\begin{equation}
r \ \leq\ \frac{E(0) + \phi_{\mathrm{cog}}(0)/\lambda(\nu) - A}{\CCC}
\label{eq:rmax-cog}
\end{equation}
を得る。\textbf{式\eqref{eq:rmax-cog}は、取引頻度が到達できる事業規模の上限を決めることを示す。}
式\eqref{eq:decay}は半減期の議論にとどまらず、第\ref{sec:growth}章の制約に接続する。
係数の値には根拠がなく、比の構造のみを示す例示である。
日常的に購入される財では認知余剰がほとんど残らず、
年単位の契約には長く残るという序列が導かれる。
\end{example}

\section{レイヤー構造：可変性による五層の整理}
\label{sec:layers}

以上を、可変性の順に五層として整理する。上の層ほど外部要因で書き換えられやすい。

\begin{center}
\begin{tabularx}{\textwidth}{@{}llX@{}}
\toprule
層 & 対象 & 可変性 \\
\midrule
制度層 & 規制、商慣行、UI規範 & 外生的に変化し、下の四層をまとめて書き換える \\
認知層 & $\dbar-\delta$ & 学習と規制で単調に削られる \\
信用層 & $\kappa_i$ の配分 & 交渉力に依存。制度層が制約を課す \\
時間層 & $\operatorname{supp}(\pi)$ と $\operatorname{supp}(\delta)$ の関係 & 物理層にかなり規定される \\
物理層 & $C_{\mathrm{prod}}$、限界費用、資産の重さ & 技術と物理法則。短期には所与 \\
\bottomrule
\end{tabularx}
\captionof{table}{可変性の順に整理した五層構造。}
\label{tab:five-layers}
\end{center}

この順序は、余剰の耐久性の順序と一致する。
認知層と制度層に依存した $\phi$ は、法改正と顧客の学習という二方向から継続的に削られる。
物理層に由来する優位は削られにくい。

\begin{remark}[制度層の作用は一様でない]
\label{rem:institution-modes}
「制度層が制約を課す」と述べたが、その形式は少なくとも三つに分かれる。

\begin{center}
\begin{tabularx}{\textwidth}{@{}lXl@{}}
\toprule
形式 & 作用 & 本稿での例 \\
\midrule
参入要件 & $\Phi$ の実行可能性を奪う & 供託義務、登録制（第\ref{sec:bootstrap}節） \\
価格規定 & 選択肢を形骸化させる & 上限制手数料（第\ref{part:results}部） \\
競争条件の操作 & $\phi_{\mathrm{barg}}$ を間接的に削る & 訴訟圧力、手数料自由化（注\ref{rem:institution-shift}） \\
\bottomrule
\end{tabularx}
\captionof{table}{制度層が制約を課す三つの作用形式。}
\label{tab:institution-modes}
\end{center}

三者は削られ方が異なる。第一と第二は特定の $\Phi$ を選べなくするのに対し、
第三は $\Phi$ を残したまま価格を下げる。
\textbf{近年の主流は第三の形式であり、価格を直接規定する形式は減少している}
（第\ref{sec:institution-mode}節）。

したがって「制度層に削られる」という記述は、
削られる速度と確実性については何も述べていない。
第\ref{sec:platform-fee}章では支配力に由来する手数料と
機能に由来する手数料が分離できる事例を扱うが、
そこで観測されるのは第三の形式による一度の引き下げであり、
継続的な減衰を示すものではない。
\end{remark}
